\documentclass[11pt]{article}
\usepackage[T1]{fontenc}
\usepackage{lmodern,microtype,amsmath,amssymb,amsthm,mathtools}
\usepackage[margin=1in]{geometry}
\usepackage{booktabs,longtable,array,enumitem}
\usepackage[dvipsnames]{xcolor}
\usepackage[colorlinks=true,linkcolor=MidnightBlue,citecolor=MidnightBlue,urlcolor=MidnightBlue]{hyperref}
\usepackage{bookmark}
\newtheorem{theorem}{Theorem}[section]
\newtheorem{proposition}[theorem]{Proposition}
\newtheorem{lemma}[theorem]{Lemma}
\newtheorem{corollary}[theorem]{Corollary}
\theoremstyle{definition}\newtheorem{definition}[theorem]{Definition}
\theoremstyle{remark}\newtheorem{remark}[theorem]{Remark}
\newcommand{\R}{\mathbb R}
\newcommand{\C}{\mathbb C}
\newcommand{\N}{\mathbb N}
\newcommand{\dd}{\,\mathrm d}
\newcommand{\ii}{\mathrm i}
\newcommand{\Res}{\operatorname{Res}}
\newcommand{\Log}{\operatorname{Log}}
\newcommand{\sech}{\operatorname{sech}}
\newcommand{\Cov}{\operatorname{Cov}}
\newcommand{\A}{\mathcal A}
\newcommand{\Eo}{\mathcal E_{\mathrm o}}
\newcommand{\Ns}{\mathcal N}
\newcommand{\Ls}{\mathcal L}
\newcommand{\Sc}{\mathcal S}
\newcommand{\Mc}{\mathcal M}
\newcommand{\Dt}{\mathcal D}
\newcommand{\abs}[1]{\lvert#1\rvert}
\newcommand{\norm}[1]{\lVert#1\rVert}
\providecommand{\researchroot}{/Users/ebbaker/Documents/shifted-zeta-positivity}
\newcommand{\src}[2]{\href{run:\researchroot/papers/investigations-5/01_heat_density_asymptotics/#1}{#2}}
\hypersetup{pdftitle={The theta heat flow: hierarchies, asymptotics, and all-time approximation},pdfauthor={Edward Baker},pdfsubject={The Book manuscript; research synthesis with explicit proof dependencies}}
\numberwithin{equation}{section}
\allowdisplaybreaks[2]
\errorcontextlines=80
\setlength{\emergencystretch}{3em}
\setlist{itemsep=.25em,topsep=.4em}
\title{The theta heat flow:\\hierarchies, asymptotics, and\\all-time approximation\\[.7em]\large The Book manuscript}
\author{Edward Baker\thanks{Drafted for Edward Baker with substantial assistance from OpenAI language models. This synthesis and its research programme were prepared with gpt-6-astra, actual recorded reasoning effort max, on 26 October 2026. Companion integration, the six-return update and proof-link reconciliation were completed with gpt-6.1-sol, actual recorded effort ultra, on 26 October 2026. The underlying notes and companion manuscripts retain their own model/effort attributions, principally gpt-6.1-sol at effort ultra and gpt-6-astra at the recorded efforts. Internal LLM assessments do not constitute independent specialist validation.}}
\date{26 October 2026\\[.25em]\small Analytical working manuscript}
\begin{document}
\maketitle
\begin{abstract}
We organize the canonical Riemann theta heat flow around three finite
approximation mechanisms: source approximation, frozen Gaussian evolution,
and expansion at a moving gamma saddle. They describe different regions
of the same entire field and admit explicit error accounts. The completed
negative hierarchy gives $O(X^{-p-3/4})$ in its gamma norm at fixed order;
controlled order growth delivers $e^{-L^2/4}$ in the outer $(3,1)$ norm
on $L^{-1}\le -t\le2$, with the distinct smaller-onset theorem on
$1\le -t\le2$. Here $L=\log(X/(4\pi))$. A retained-arithmetic model
reaches zero from the positive side and, with corrected two-arm and
compact saddle charts, gives a high-height field atlas on $[0,T_*]$
for each fixed finite $T_*\ge33$. Every fixed positive saddle order
also has an eventual all-height theorem. Near the prescribed regulator
switch, fixed $P_4$ gives the full physical tuple: static potential and
density allowances $X^{-7}$, current/covariance $X^{-7}(1+\log X)^2$,
and specified time/source outputs $X^{-6}(1+\log X)^2$, with retained
heat residual corrections and output-specific derivative orders.
These regional results stand beside the existing finite all-time,
all-height regulated physical construction, with actual covariance and
all joining defects. Source compression and finite evaluation have
polynomial additional-accuracy cost at fixed data and query box.
Explicit constants/onsets, lower margins, isolated-cluster initial data,
complete exterior/energy payments, bounded-height refinements, stronger
negative precision through zero, new physical consumers and useful joint
cost remain distinct tasks. We present the
principal formulas and proof mechanisms, identify supporting proofs, and
state the exact scope of each conclusion.
\end{abstract}
\paragraph{Why ``Book''.}
The name follows the author's reference to Erd\H{o}s's ideal of proofs
from \emph{The Book}: economical formulas, visible mechanisms and errors
that can be understood and improved. It is an editorial and research
aspiration. This manuscript is a synthesis of the present results, with
its unfinished bridges displayed; it does not claim a completed best
possible description or a proof of the Riemann hypothesis.
\clearpage
{\small\tableofcontents}
\clearpage
\section{One field, three approximation mechanisms}
\label{sec:canonical}

\subsection{The canonical object}
For $u\ge0$ set
\begin{equation}
 \Phi(u)=\sum_{n\ge1}
 (2\pi^2n^4e^{9u}-3\pi n^2e^{5u})e^{-\pi n^2e^{4u}},
 \qquad H(t,z)=\int_0^\infty e^{tu^2}\Phi(u)\cos(zu)\dd u.
 \label{eq:source}
\end{equation}
The source has a smooth even extension, is positive on the real axis,
and decays faster than every Gaussian. Its full arithmetic series is
part of the definition. Termwise differentiation on compact sets gives
\begin{equation}
 H_t=-H_{zz},\qquad
 H(0,z)=\frac18\xi\left(\frac{1+\ii z}{2}\right),\qquad
 \frac18\xi(w)=M(w)\zeta(w).
 \label{eq:normalization}
\end{equation}
The gamma factor is
\[
 M(w)=\frac{w(w-1)}{16}\pi^{-w/2}\Gamma(w/2).
 \]
Thus $H$ is entire in $z$, real on the real axis, and even. We call
$x=\Re z$ the height and $y=\Im z$ the transverse coordinate.
Throughout, ``all heights'' means every real $x$ on a specified bounded
transverse strip, not every complex $z$ with a strip-independent constant.

Two exact evolution formulas explain the different useful approximations.
For $T=-t>0$ and $s=(1+\ii z)/2$,
\begin{align}
 H(-T,z)&=\frac1{\sqrt{4\pi T}}
   \int_\R e^{-(z-u)^2/(4T)}H(0,u)\dd u,\label{eq:negative-heat}\\
 H(t,z)&=\frac1{\sqrt{\pi t}}
   \int_\R e^{-v^2/t}\frac{\xi(s+v)}8\dd v\quad(t>0).
 \label{eq:positive-heat}
\end{align}
The first evolves the initial field by forward heat flow in $T$.
The second moves the gamma saddle into a right half-plane, where the
ordinary Dirichlet series can be estimated in full. These identities
use no hypothesis about zeros. The normalization and canonical contour
input are recorded in the parent manuscript and its references
\cite{parent,positive,poly}.

\subsection{What is already global, and what remains regional}
There are three distinct achievements. First, for every prescribed
tolerance there is one finite construction approximating a specified
regulated, averaged physical tuple at every real time and height
(Theorem~\ref{thm:global}). Its constants are large but finite, its
covariance is the actual covariance, and its joining errors are paid.
Second, there is a high-height field atlas reaching zero time. Third, a specialized fixed-$P_4$ consumer proves a sharper physical tuple on $[63,65]$. These have different norms and quantifiers. There are also sharper asymptotic formulas on other particular regions.
Table~\ref{tab:atlas} displays their leading field allowances. It is an
atlas of proved domains, not a single theorem permitting free exchange
of their quantifiers.

\begin{table}[ht]
\centering\small
\begin{tabular}{p{.22\textwidth}p{.30\textwidth}p{.38\textwidth}}
\toprule
Region & Approximation & Allowance and scope\\
\midrule
$t=-T\to-\infty$ & Gaussian endpoint & Absolute $O(T^{-3/2})$ field error on $|y|\le A\sqrt T$, all $x$.\\
$\tau\le T\le B$ & Completed negative hierarchy, order $p$ & $O(X^{-p-3/4})$ in the gamma norm; fixed $p,\tau,B$.\\
$1\le T\le2$ & Order $\lfloor\tfrac12\log(X/(4\pi))\rfloor$ & $e^{-\log^2(X/(4\pi))/4}$ in the outer norm.\\
$L^{-1}\le T\le2$ & Integrated moving-endpoint hierarchy & $e^{-L^2/4}$ outer $(3,1)$ at $L\ge1024\nu^2$; strict query/endpoint interfaces.\\
$0\le t\le1$ & Frozen retained-arithmetic $F_N$ & $\epsilon_X=\Ls^2X^{-3/4}$ in $\Eo$; stronger $X^{-3/4}$ near zero.\\
$3/4\le t\le33$ & Corrected elementary $F_1$ & $B_1=\Ls^4/X^2+X^{-t\log2/4}$ in $\Eo$.\\
$|t|\log(X/(4\pi))\le8$ & Completed crossover carrier & $X^{-3/4}e^{t_-\log^2(X/(4\pi))/16}$ in the outer envelope.\\
$t\ge t_{*,J}$ & Positive saddle $P_{2J}$ & $t^{-J-1}\log^{-2J-3}(2+t)$, every height, each fixed $J$.\\
Fixed $[a,T]\subset[32,\infty)$ & Positive saddle $P_{2J}$ & $X^{-2J-3}+X^{-t\log2/4}$; constants depend on $J,a,T$.\\
$63\le t\le65$, high $X$ & Fixed $P_4$ physical consumer & Absolute output-specific rates in Theorem~\ref{thm:p4-physical}; fixed regulator/filter, no global join.\\
Every $t,x$, $|y|\le1$ & Finite global assembly & Arbitrary prescribed absolute tolerance for the full physical tuple.\\
\bottomrule
\end{tabular}
\caption{Principal established regimes. Spatial and time losses, residuals,
buffers and effectivity qualifications are given in the corresponding sections.}
\label{tab:atlas}
\end{table}

The global construction is therefore a foundation, not a reason to stop
improving the atlas. The research aim is to replace costly intermediate
pieces with transparent, refinable formulas and to transfer their sharper
errors through the nonlinear physical observables without losing the gain.

\subsection{A shared finite algebra}
Both hierarchies start from a finite weighted expansion of an exponential.
If $c_r$ are coefficients and $d$ is fixed, define
\begin{equation}
 A_\ell(q)=\sum_{\substack{\nu_1,\ldots,\nu_\ell\ge0\\
                          \sum r\nu_r=\ell}}
 q^{\ell+d|\nu|}\prod_{r=1}^{\ell}\frac{c_r^{\nu_r}}{\nu_r!},
 \qquad A_0=1.
 \label{eq:finite-partition}
\end{equation}
Finite coefficient extraction yields
\begin{equation}
 \ell A_\ell(q)=\sum_{r=1}^{\ell}r c_r q^{r+d}A_{\ell-r}(q).
 \label{eq:shared-recurrence}
\end{equation}
For negative time $d=1$: the expansion retains curvature beyond a
linear gamma carrier. For positive time $d=2$: the saddle and its
quadratic Gaussian have already been removed, so odd grades integrate
to zero. The algebra is common; the analytic remainder and its domain
must be proved separately. A formal recurrence alone implies neither
convergence nor a permissible order growing with a large parameter.

\subsection{Conventions that preserve the meaning of an error}
We use $L=\log(X/(4\pi))$ and $\Ls=1+\log X$ on high-height charts, and reserve $\Lambda$ for the Newman constant. The large-time saddle section uses its explicitly named scale $\mathcal L_+=\log R$. A chart centered at $X$ comes with a larger analytic rectangle.
Spatial derivatives are estimated only after that larger rectangle is
paid. The center $X$, arithmetic cutoff, hierarchy order and real
integration endpoints are frozen under differentiation. A time or
height dependent choice of order describes a family of formulas; it is
not silently differentiated as one smooth formula.

Field errors are divided by a positive sum of arm moduli, never by
$H$ or by a signed approximation at its zeros. We keep four accounts:
field and spatial jets; direct time error; the exact model residual;
and evaluation error. Nonlinear physical transfer introduces a fifth
account involving the prescribed regulator and averaging kernel.
None of these accounts can replace another without an explicit identity.

The results quoted here retain the status of their supporting notes.
This manuscript reorganizes their proofs and principal formulas; it is
not an independent specialist verification of every source constant.
Appendix~\ref{sec:sources} specifies the proof dependencies and the
division of work among the supporting manuscripts.

\section{Negative time: smoothing and the completed fluctuation hierarchy}
\label{sec:negative}

\subsection{The remote negative endpoint}
Let $\phi_0=\Phi(0)>0$, $C=\phi_0\sqrt\pi/2$, and
\begin{equation}
 N_T(z)=C T^{-1/2}e^{-z^2/(4T)}.
 \label{eq:negative-gaussian}
\end{equation}
An explicit convergent source-series majorant supplies
$B_{2,\Phi}\ge\sup_{u\ge0}|\Phi''(u)|$. Evenness gives
$|\Phi(u)-\phi_0|\le B_{2,\Phi}u^2/2$. Substitution in
\eqref{eq:source}, followed by $v=\sqrt T u$, proves
\begin{equation}
 |\partial_z^j(H(-T,z)-N_T(z))|
 \le\frac{B_{2,\Phi}}2 T^{-(j+3)/2}J_{j+2}(A),
 \quad |y|\le A\sqrt T,
 \label{eq:negative-absolute}
\end{equation}
for each fixed derivative order $j$, uniformly in real $x$. Here
\begin{align*}
 J_m(A)&=\int_0^\infty v^m e^{-v^2+Av}\dd v,\\
 J_0(A)&=\frac{\sqrt\pi}2e^{A^2/4}\operatorname{erfc}(-A/2),
 &J_1(A)&=\frac{1+AJ_0(A)}2,\\
 J_{m+1}(A)&=\frac{mJ_{m-1}(A)+AJ_m(A)}2\quad(m\ge1).
\end{align*}
This elementary proof also explains the normalization of the Gaussian.
It is an absolute estimate. At very large heights the Gaussian can be
far smaller than its error; no uniform relative high-height asymptotic
is asserted.

For a fixed constant regulator $\epsilon>0$, the unaveraged density
defined in Section~\ref{sec:physical} has the scaled endpoint profile
\begin{equation}
 T^2\rho_{\epsilon}(-T,\sqrt T(\xi+\ii\upsilon))
 =\frac{\phi_0^2}{16\epsilon^2}(\xi^2+\upsilon^2)
       e^{-(\xi^2-\upsilon^2)/2}+O(T^{-1})
 \label{eq:negative-density-profile}
\end{equation}
on compact $(\xi,\upsilon)$ sets. The physical adaptive regulator used
later is different; its endpoint transfer has its own error rates.
The parent endpoint proof gives both accounts \cite{parent}.

\subsection{Frozen high-height charts}
Fix $0<\tau\le T\le B$, observation widths $r,b$, and a large real
center $X$. Write
\[
 L=\log(X/(4\pi)),\quad N=\lfloor\sqrt{X/(4\pi)}\rfloor,
 \quad V_B=8\sqrt{\max(1,B)}L,\quad d=z-X.
\]
All of these choices are frozen on the chart. Put
\[
 w_\epsilon(z)=\frac{1+\epsilon\ii z}2,\quad
 \ell_\epsilon(z)=\log M(w_\epsilon(z)),\quad
 G(z)=\sqrt{M(w_+(z))M(w_-(z))},\qquad \epsilon\in\{+1,-1\},
\]
with compatible analytic branches and $G(X)>0$. On fixed buffers,
$|G(z)|$ has the scale $X^{7/4}e^{-\pi X/8}$ up to bounded ratios.
The gamma norm is
\begin{equation}
 \|E\|_\gamma^{(3,1)}=
 \sup_{\substack{\tau\le T\le B\\|\Re z-X|\le r,\ |\Im z|\le b}}
 \frac{\max_{0\le j\le3}|\partial_z^j E|+|\partial_T E|}{|G(z)|}.
 \label{eq:gamma-norm}
\end{equation}

For each arm and arithmetic mode set
\begin{align*}
 b_{\epsilon n}&=M(w_\epsilon(X))n^{-w_\epsilon(X)},\\
 \lambda_{\epsilon n}&=\ell_\epsilon'(X)-\epsilon\ii\log n/2,
 &c_{\epsilon r}&=\frac{\ell_\epsilon^{(r+1)}(X)}{(r+1)!}.
\end{align*}
The fluctuation factor is
\[
 \exp\{\ell_\epsilon(X+q)-\ell_\epsilon(X)-\ell_\epsilon'(X)q\}.
\]
Its finite weighted polynomials are \eqref{eq:finite-partition} with
$d=1$ and $c_r=c_{\epsilon r}$. A simultaneous Cauchy estimate on a
large disk gives $|c_{\epsilon r}|\le A D^{r-1}X^{-r}$ for every needed
index. This common bound, rather than separate unspecified constants
at each order, makes controlled order growth possible.

Gaussian evolution of a polynomial is explicit. With
\begin{equation}
 \mathcal H_m(u,T)=m!\sum_{\nu=0}^{\lfloor m/2\rfloor}
                 \frac{T^\nu u^{m-2\nu}}{\nu!(m-2\nu)!},
 \quad \mathcal H_m=u\mathcal H_{m-1}+2T(m-1)\mathcal H_{m-2},
 \label{eq:heat-moments}
\end{equation}
the heat transform of $q^m e^{\lambda q}$ is
$e^{\lambda d+T\lambda^2}\mathcal H_m(d+2T\lambda,T)$.
Equivalently, define $Q_T=d+2T\partial_d$ and
\begin{align}
 U_{\epsilon0}&=\sum_{n\le N}b_{\epsilon n}
                          e^{\lambda_{\epsilon n}d+T\lambda_{\epsilon n}^2},\\
 jU_{\epsilon j}&=\sum_{r=1}^j r c_{\epsilon r}
                            Q_T^{r+1}U_{\epsilon,j-r}.
 \label{eq:negative-recurrence}
\end{align}
Thus every retained grade is a finite Gaussian evolution of fixed real
initial data and solves the heat equation exactly in $T$.

\subsection{The arithmetic completion is part of the model}
A bare frozen sum misses cutoff crossings and the Riemann--Siegel
remainder. On $[X-V_B,X+V_B]$ let
\[
 a(u)=\sqrt{u/(4\pi)},\quad m(u)=\lfloor a(u)\rfloor,
 \quad\chi(u)=1-2(a(u)-m(u)),
\]
and let $S_m(u)=\sum_{\epsilon,n\le m}M(w_\epsilon(u))n^{-w_\epsilon(u)}$.
Use the fixed real Riemann--Siegel correction
\begin{align*}
 \mathcal U(u)&=\exp\{-\ii[u\log(u/(4\pi))/4-u/4-\pi/8]\},\\
 R_{m,K}(u)&=2\Re\left\{(-1)^{m-1}M(w_+(u))\mathcal U(u)a(u)^{-1/2}
       \sum_{k=0}^K C_k^{\rm RS}(\chi(u),1/2)a(u)^{-k}\right\}.
\end{align*}
The coefficient functions $C_k^{\rm RS}$ are the fixed canonical
functions in the supporting recurrence manuscript, with its finite
providers and removable-value conventions. Define
\begin{align}
 I_{X,K,B}(T,z)&=\int_{X-V_B}^{X+V_B}
 \frac{e^{-(z-u)^2/(4T)}}{\sqrt{4\pi T}}
       [S_{m(u)}(u)-S_N(u)+R_{m(u),K}(u)]\dd u,\\
 F^{[p,K]}(T,z)&=\sum_{\epsilon=\pm1}\sum_{j=0}^p U_{\epsilon j}(T,z)
                     +I_{X,K,B}(T,z).
 \label{eq:negative-completed}
\end{align}
The integrand is evaluated as fixed real data on its fixed partition.
No moving arithmetic floor is differentiated. Consequently
$(\partial_T-\partial_{zz})F^{[p,K]}=0$ exactly. Finite-precision outputs
have a separate evaluation error; their rounded values are not themselves
claimed to satisfy an exact differential equation.

\begin{theorem}[Completed fixed order; recurrence manuscript]
\label{thm:negative-fixed}
For fixed $p\ge1$, $0<\tau\le B$ and fixed observation buffers, there
are finite constants and a sufficient height onset such that
\begin{equation}
 \|H(-T,\cdot)-F^{[p,2p]}(T,\cdot)\|_\gamma^{(3,1)}
          \le C_{p,r,b}(\tau,B)X^{-p-3/4}.
 \label{eq:negative-fixed-error}
\end{equation}
\end{theorem}
\begin{proof}[Proof mechanism]
The local weighted-exponential remainder is paid with its Gaussian
moments. Summing the absolute arithmetic population costs $X^{1/4}$.
The canonical Riemann--Siegel initial remainder is bounded by
\[
 \frac{2^{7/4}}7G(u)a(u)^{-1/2}
       \Gamma((K+1)/2)\left(\frac{1.1}{a(u)}\right)^{K+1},
\]
so $K=2p$ gives the same power $X^{-p-3/4}$ after heat propagation.
The cutoff correction pays every jump; the omitted real heat tails
are kept separately. Spatial and time kernel derivatives are bounded
on larger buffers before restriction. This proves the asserted norm
with the explicit finite constant construction in \cite{recurrence}.
\end{proof}
The detailed proof is the recurrence companion's
\src{recurrence/manuscript/sections/04_comparison.tex}{fixed-order theorem}
(\texttt{thm:fixed}). The constants contain inverse powers of $\tau$ and factors such as
$e^{(b+1)^2/(4\tau)+2B}$. The theorem does not allow an unproved
substitution $\tau=\tau(X)$ or $B=B(X)$.

\subsection{Growing order, onset and finite evaluation}
Passing from the gamma envelope to the outer envelope of
Section~\ref{sec:crossover} costs at most a fixed multiple of
$e^{L^2/8}$ on $1\le T\le2$. The simultaneous coefficient estimates
control the remaining order dependence by $e^{O(L\log L)}$.
Keeping the heat tails leads to the following stronger conclusion
\cite{recurrence,rec-onset,rec-bounds}.

\begin{theorem}[Controlled growing order]
\label{thm:negative-growing}
For fixed $a>1/8$, $p=\lfloor aL\rfloor$, and
$0<\eta<\min(a-1/8,127/8)$, the completed error on $1\le T\le2$
is at most $e^{-\eta L^2}$ in the outer version of
\eqref{eq:gamma-norm}, beyond a sufficient onset. In particular
$p=\lfloor L/2\rfloor$ permits $\eta=1/4$.
\end{theorem}

The recurrence companion gives the full growing-order proof
(\src{recurrence/manuscript/sections/05_growing_order.tex}{\texttt{thm:growing}}).
For this specified $a=1/2$ construction, write
$\nu=\lceil r+b+2\rceil+1$. The sharpened sufficient analytic onset is
\begin{equation}
 L\ge64\nu;
 \qquad X\ge2^{128\nu+8}\quad\hbox{is a sufficient dyadic choice}.
 \label{eq:rec-onset}
\end{equation}
The older $L\ge2^{22}\nu^4$ onset remains a valid fallback, not the
preferred requirement. Strict domain and denominator guards remain
part of the certified descriptors.

The finite evaluator returns the field, spatial derivatives through
three and the time derivative with a combined analytic and evaluation
allowance $\Eo e^{-L^2/4}$. For rational query data and dyadic centers,
the delivered precision has scale
\[
 P=\frac{3L^2}{8\log2}+O(L\log L).
\]
With $Q$ the sum of input length, delivered precision, separation bits,
$(\lceil L\rceil)^2$ and a fixed offset, the proved coarse bounds are
\begin{equation}
 \operatorname{work}\le C(N+1)Q^{40},\qquad
 \operatorname{memory}\le C Q^{20}.
 \label{eq:negative-complexity}
\end{equation}
The arithmetic population $N\asymp\sqrt X$ has not disappeared.
The uniform majorants $B_F$ and every $D_k$ are multiplied by
$\operatorname{csch}(\pi)<1/10$. Separately, the enlarged-contour
majorant $B_c/16$ permits eleven fewer Gauss nodes per edge,
forty-four per instantiated four-edge rule, at the same target.
These are analytical contour-count improvements; they do not reduce
the outer completion-integral nodes or establish measured runtime gains.
The exact input guards, half-target onset and delivered evaluation contract
are in the recurrence companion's
\src{recurrence/manuscript/sections/11_calibration_and_onset.tex}{\texttt{thm:linear-onset}}
and \src{recurrence/manuscript/sections/12_effective_evaluation.tex}{\texttt{thm:effective}, \texttt{thm:bit-cost}}. Component replays support the finite provider recipes;
a complete software implementation is a separate task.

\subsection{The integrated moving lower endpoint}
Recurrence Note 20's proof is integrated as
\src{recurrence/manuscript/sections/14_moving_endpoint.tex}{\texttt{thm:moving-endpoint-delivery}},
with its distinct \texttt{def:moving-channel-norm} \cite{rec-moving}.
It delivers field, spatial orders zero through three and direct time
error together in the outer $(3,1)$ norm on $L^{-1}\le T\le2$:
\[
 \|H(-T,\cdot)-\widehat F\|_{\mathrm o,[1/L,2]}^{(3,1)}
 \le e^{-L^2/4},\qquad L\ge1024\nu^2.
\]
The construction keeps $p=\lfloor L/2\rfloor$, $K=2p$,
$V=8\sqrt2L$, all arithmetic modes, exact adjacent-cutoff transfers
and the bounded completion integral. The new proof pays
$e^{Y^2/(4T)}$, the Gaussian normalization, inverse-time powers,
complete population and degree-dependent evolved moments jointly.
Substitution into a fixed-$\tau$ theorem would not supply these payments.

Ordinary rational queries carry a strict witness $T L^->1$ for a
certified rational lower bound $L^-\le L$. The literal $T=1/L$ tag
has a separately paid time-input error channel, including $F_{TT}$.
There is no arbitrary-real equality oracle. Delivered precision remains
$3L^2/(8\log2)+O(L\log L)$, but the unchanged outer ellipse now needs
$O_{r,b}(L^3)$ sufficient completion-quadrature nodes. Internal provider
work and query separation are additional costs. The preserved $[1,2]$
theorem retains its smaller $64\nu$ onset and its own outer-node count.
Neither theorem includes $T=0$ or a complete certified software evaluator.

\section{The height carrier and the narrow crossover through zero}
\label{sec:crossover}

On a right-hand high-height chart use $L_z=\Log(z/(4\pi))$, continued
from the positive real axis, and define
\begin{align}
 A_o(t,z)&=\frac{\pi^{1/4}z^{7/4}}{16}
  \exp\left\{-\frac{\pi z}8+\frac t{16}(L_z^2-\pi^2/4)\right\},\\
 \Theta_o(t,z)&=\frac z4(L_z-1)+\frac{\pi tL_z}{16}+\frac{7\pi}8,
 \quad C_\pm=\tfrac12 A_o e^{\pm\ii\Theta_o},
 \quad \Eo=|C_+|+|C_-|.
 \label{eq:outer-arms}
\end{align}
The phase is unwrapped as written; each arm has half the cosine
amplitude. The envelope includes transverse growth. A real-axis
amplitude alone would not pay the buffered comparison.

Put $a(z)=\sqrt{z/(4\pi)}$, freeze $N=\lfloor a(X)\rfloor$, and set
\begin{align*}
 w_n(t,z)&=n^{-1/2}\exp\{t[(\log n)^2-L_z\log n]/4\},\\
 \varphi_n(t,z)&=\Theta_o(t,z)-(z/2+\pi t/8)\log n,\\
 \mathfrak C(v)&=\frac{\cos(2\pi(v^2-v)-\pi/8)}{\cos(2\pi v)},\\
 R_N(z)&=(-1)^N G(z)a(z)^{-1/2}\mathfrak C(a(z)-N),\\
 F_N(t,z)&=A_o(t,z)\sum_{n\le N}w_n(t,z)\cos\varphi_n(t,z)
                           +e^{-\pi^2t/64}R_N(z).
\end{align*}
Values at the apparent poles of $\mathfrak C$ are taken by removable
continuation. The identity
\begin{equation}
 \mathfrak C(v)+\mathfrak C(v-1)
        =2\cos(2\pi(v-1)^2-\pi/8)
 \label{eq:cutoff-identity}
\end{equation}
pays the adjacent cutoff change. This leading completion is simpler
than the high-order completion of Section~\ref{sec:negative}; it also
has a weaker error.

\begin{theorem}[Buffered signed crossover]
\label{thm:crossover}
On fixed buffered charts with $L_X=\log(X/(4\pi))$ and
$|t|L_X\le8$, let $t_-=\max(0,-t)$ and
$\delta=X^{-3/4}e^{t_-L_X^2/16}$. The completed canonical estimate is
\begin{align}
 |\partial_z^j(H-F_N)|&\le C_j\Eo\delta(1+L_X)^j,
                         &&0\le j\le3,\\
 |\partial_t(H-F_N)|&\le C\Eo\delta(1+L_X)^2,\\
 |(\partial_t+\partial_{zz})F_N|&\le C\Eo\delta(1+L_X).
 \label{eq:crossover-errors}
\end{align}
In the whole signed window $\delta=O(X^{-1/4})$; on its positive
half $\delta=X^{-3/4}$.
\end{theorem}
The proof retains both canonical arms, the complete arithmetic
population and cutoff payment on the larger strip. The time estimate
comes from the exact heat error identity and the actual residual; it
does not introduce a singular $1/t$ by differentiating an error
allowance. See \src{ansatz/notes/04_COMPLEX_STRIP_REMAINDER_AND_BUFFERED_CROSSOVER_20261025.md}{Ansatz Note 04}
and the supporting manuscript's
\src{ansatz/manuscript/main.tex}{buffered crossover estimates}
(\texttt{eq:buffered-crossover-jets}, \texttt{eq:buffered-crossover-time})
\cite{ansatz}.

This theorem explains a genuine but narrow two-sided passage through
zero. It does not extend the $e^{-L^2/4}$ recurrence estimate to zero.
An earlier paid chartwise negative-time join on $[-2,0]$ also uses
its own weaker representation and error account
(\src{ansatz/notes/06_NEGATIVE_BRIDGE_AND_TIME_JOIN_20261025.md}{Ansatz Note 06}).

\subsection{Retained arithmetic reaches the positive compact band}
The same frozen $F_N$, including its explicit completion, now has the
longer-band contract of Ansatz Note 26 \cite{retained}:
\begin{align}
 |\partial_z^j(H-F_N)|&\le C_j\Eo\Ls^j\epsilon_X,
       &&j\le3,\quad 4/L\le t\le1,\notag\\
 |(H-F_N)_t|&\le C_t\Eo\Ls^2\epsilon_X,
       &\epsilon_X&=\Ls^2X^{-3/4},\label{eq:retained-band}\\
 |\Res F_N|&\le C\Eo\epsilon_X,
 &|\partial_z\Res F_N|&\le C\Eo\Ls\epsilon_X.\notag
\end{align}
There are $2N=O(\sqrt X)$ arms. The proof supplies a growing-shift
initial-error lemma for the positive Gaussian's imaginary shifts
$O(tL)$. It uses a fixed inner strip and, outside it, the complete
ordinary Dirichlet tail, gamma ratios and floor-sensitive powers.
The full mode population costs $X^{1/4}$, while curvature/calibration
costs $\Ls^2/X$. Integrated Gaussian moments pay direct time before
using $1/t\le L/4$. These mechanisms explain the stated allowance;
a fixed-strip theorem alone would not cover the new shifts.

On $4\le tL\le8$ this is exactly the crossover formula, including
branches, normalization and completion. Its model difference is zero.
Consequently the same $F_N$ covers $0\le t\le1$ at the coarser
$\epsilon_X$ allowance, retaining the stronger inherited $X^{-3/4}$
allowance on $0\le tL\le8$. No time splice or derivative of $N$ is used.
The full proof and union are in the ansatz companion's
\src{ansatz/manuscript/08_positive_bridges.tex}{\texttt{thm:retained-positive-band}, \texttt{thm:positive-atlas}}.

Ansatz Note 24 separately proves $F_0=C_++C_-$ on $[1/2,1]$:
field/spatial $j\le3$ and direct-time allowances are respectively
$C_j\Eo\Ls^jX^{-t\log2/4}$ and
$C_t\Eo\Ls^2X^{-t\log2/4}$. The actual residual and its spatial
derivative are $C\Eo\Ls^2/X$ and $C\Eo\Ls^3/X$.
Its full-buffer calibration at $t=1/2$ is proved anew; a restricted-strip
reference theorem is not extended by citation \cite{small-positive}.
The ansatz \texttt{thm:small-positive-slab} contains the detailed proof.

Together with $F_1$ on $[3/4,33]$ and the compact saddle theorem on
$[32,T_*]$, these give a buffered high-height field atlas on $[0,T_*]$
for every fixed finite $T_*\ge33$, after the maximum of the actual
$T_*$-dependent onsets. Section~\ref{sec:positive} states the genuine
$F_1$ overlaps. This is an atlas of fields, with paid comparisons through
$H$, rather than one newly assembled smooth global field. Bounded-height
coverage, sharper physical consumers and useful joint cost are separate.
The stronger negative $e^{-L^2/4}$ precision still excludes zero; its
possible extension to $T\asymp L^{-2}$ is not supplied by this union.

\section{Positive time: the moving saddle and every fixed correction order}
\label{sec:positive}

\subsection{The saddle absorbs the large motion}
Write $\alpha=(\log M)'$. For each $w=s,1-s$, use the large
right-half-plane solution
\begin{equation}
 p=w+\frac t2\alpha(p),\qquad
 B=1-\frac t2\alpha'(p),\qquad
 \Sc_t(w)=M(p)e^{-(p-w)^2/t}B^{-1/2}.
 \label{eq:positive-saddle}
\end{equation}
The logarithm of $M$ is real on $p>1$ and $B^{-1/2}$ is the branch
near one. Fix the observation widths and a positive analytic margin;
all estimates start on the rectangle with twice that margin. For a
real center $x$ put
\begin{equation}
 X=|x|,\quad R=2+t+X,\quad \mathcal L_+=\log R,\quad
 \Ns=1+X+t\mathcal L_+,\qquad
 \A=|\Sc_t(s)|+|\Sc_t(1-s)|.
 \label{eq:positive-scales}
\end{equation}
At a sufficient onset, contraction in a fixed sector gives
\begin{equation}
 |p|\asymp\Ns,\quad \Re(p-w)\ge t\mathcal L_+/5,\quad
 |\Im(p-w)|\le t,\quad |B-1|\le1/4,\quad |p|\ge t.
 \label{eq:saddle-gates}
\end{equation}
These gates include bounded heights and the real center $x=0$.

\subsection{Finite coefficients}
Let $a_r=\alpha^{(r)}(p)$, $\beta_r=a_{r+1}/(r+2)!$, and
$\mathcal P_\ell=\{\nu\ge0:\sum r\nu_r=\ell\}$. With $k=|\nu|$,
define
\begin{align}
 A_\ell(u)&=\sum_{\nu\in\mathcal P_\ell}u^{\ell+2k}
                       \prod_r\frac{\beta_r^{\nu_r}}{\nu_r!},
 &\ell A_\ell&=\sum_{r=1}^\ell r\beta_r u^{r+2}A_{\ell-r},\\
 C_{2j}^{+}&=\sum_{\nu\in\mathcal P_{2j}}
 \frac{(2j+2k-1)!!}{2^{j+k}}
 \left(\frac tB\right)^{j+k}
 \prod_r\frac{a_{r+1}^{\nu_r}}{(r+2)!^{\nu_r}\nu_r!}.
 \label{eq:positive-coefficients}
\end{align}
The empty partition gives $C_0^+=1$. The superscript distinguishes
these coefficients from the negative-time Riemann--Siegel coefficients.
Parity kills every odd grade in the Gaussian moment. The first two
nonconstant coefficients are
\begin{align}
 C_2^+&=\frac{t^2\alpha'''(p)}{32B^2}
               +\frac{5t^3\alpha''(p)^2}{192B^3},\label{eq:c2}\\
 C_4^+&=\frac{t^3\alpha^{(5)}(p)}{384B^3}
 +\frac{t^4}{B^4}\left\{\frac{7\alpha''(p)\alpha^{(4)}(p)}{768}
                  +\frac{35\alpha'''(p)^2}{6144}\right\}\notag\\
 &\hspace{1em}+\frac{35t^5\alpha''(p)^2\alpha'''(p)}{2048B^5}
                  +\frac{385t^6\alpha''(p)^4}{73728B^6}.
 \label{eq:c4}
\end{align}
The approximation is
\begin{equation}
 P_{2J}(t,z)=\sum_{w=s,1-s}\Sc_t(w)\sum_{j=0}^J C_{2j}^+(t,w).
 \label{eq:positive-model}
\end{equation}
Every sum is finite. Nothing in this definition asserts convergence
of the infinite correction series.

\subsection{The complete error, before simplification}
For $n=2J+2$ define
\begin{equation}
 b_J=\sum_{k=1}^{2J+2}\frac{t^{J+1+k}}{\Ns^{2J+2+k}},\qquad
 \delta_J=b_J/t,\qquad
 E_J=b_J+R^{-\kappa_a t}+R^{b_0}e^{-ct\mathcal L_+^2},
 \quad\kappa_a=\frac{\log2}{8}.
 \label{eq:positive-budget}
\end{equation}
Here $b_0$ is a fixed middle-strip exponent and $c>0$ may be decreased
to at most $1/400$. The last two terms pay the complete arithmetic
remainder and the nonlocal contour pieces.

\begin{theorem}[Every fixed positive order]
\label{thm:positive-fixed}
Fix $J\ge0$ and fixed observation buffers. There are finite constants
and $R_{0,J}$ such that for $t\ge32$ and $R\ge R_{0,J}$,
\begin{align}
 |\partial_z^d(H-P_{2J})|&\le C_{J,d}\A\mathcal L_+^dE_J,
                         &&0\le d\le3,\\
 |(H-P_{2J})_t|&\le C_{t,J}\A(\mathcal L_+^2E_J+\delta_J),\\
 |\Res P_{2J}|&\le C_{r,J}\A\delta_J,
 &\Res&=\partial_t+\partial_{zz}.
 \label{eq:positive-full-error}
\end{align}
There is a sufficient $t_{*,J}$ for which, at every real height,
\begin{align}
 \frac{|\partial_z^d(H-P_{2J})|}{\A}
  &\le\frac{C_{J,d}}{t^{J+1}\log^{2J+3-d}(2+t)},&&0\le d\le3,\\
 \frac{|(H-P_{2J})_t|}{\A}
  &\le\frac{C_{t,J}}{t^{J+1}\log^{2J+1}(2+t)},\\
 \frac{|\Res P_{2J}|}{\A}
  &\le\frac{C_{r,J}}{t^{J+2}\log^{2J+3}(2+t)}.
 \label{eq:positive-eventual}
\end{align}
The field bound is first proved on the larger analytic rectangle;
the jet bounds use its smaller interior.
\end{theorem}
The full proof is now integrated in the ansatz companion's
\src{ansatz/manuscript/07_positive_hierarchy.tex}{\texttt{thm:positive-fixed}},
from the original Ansatz Note 19 \cite{positive}. Its
quantifier is ``for every fixed $J$, there exist constants and an
onset.'' The constants are specified by finite constructions and named
source constants, but are not numerically instantiated. The theorem
does not make $t=32$ an all-height onset for every $J$.

\subsection{Why the finite remainder has this form}
The following argument is the central reason the hierarchy is
systematically improvable at fixed order. Set $f=\log M$,
$a=\Re(p-w)$, $m=|p|$, and use the actual symmetric central segment
$\mathcal U=[-a/2,a/2]$. For $0\le\lambda\le1$ introduce
\begin{equation}
 E_\lambda(u)=u^2\int_0^1(1-r)
                  [f''(p+r\lambda u)-f''(p)]\dd r.
 \label{eq:interpolation}
\end{equation}
At $\lambda=1$ this is the exact nonquadratic exponent; at zero it
vanishes. Sectorial gamma estimates and Cauchy disks supply finite
$K_r$ with $|\alpha^{(r)}(q)|\le K_r|q|^{-r}$. After increasing the
onset, $|p+r\lambda u|\ge m/3$ and
\begin{equation}
 |\partial_\lambda^rE_\lambda|
 \le d_r\frac{|u|^{r+2}}{m^{r+1}},\qquad
 d_r=\frac{3^{r+1}K_{r+1}}{(r+1)(r+2)},\quad1\le r\le n.
 \label{eq:lambda-derivatives}
\end{equation}
The full interpolated Gaussian, not the nonquadratic factor by itself,
satisfies
\begin{equation}
 |e^{-Bu^2/t+E_\lambda(u)}|\le e^{-u^2/(2t)}.
 \label{eq:common-gaussian}
\end{equation}
Indeed its exponent is $-u^2/t+u^2\int_0^1(1-r)f''(p+r\lambda u)\dd r$;
the sufficient gate $3K_1t/(2m)\le1/2$ pays the second term.

Define the finite Bell constants
\[
 D_{n,k}=\sum_{\substack{\nu\in\mathcal P_n\\|\nu|=k}}
                   \prod_{r=1}^n\frac{(d_r/r!)^{\nu_r}}{\nu_r!}.
\]
Real Taylor's theorem in $\lambda$ through degree $n-1$ gives
\begin{equation}
 \left|e^{-Bu^2/t}\left(e^{E_1(u)}-\sum_{\ell=0}^{n-1}A_\ell(u)\right)\right|
 \le e^{-u^2/(2t)}\sum_{k=1}^nD_{n,k}
                            \frac{|u|^{n+2k}}{m^{n+k}}.
 \label{eq:bell-remainder}
\end{equation}
Every operation is finite, and only derivatives through
$\alpha^{(2J+3)}$ are needed. Integrating the even powers against
the Gaussian yields $t^{n/2+k}/m^{n+k}$, hence $b_J$.
Odd grades cancel on the actual symmetric segment. Only the retained
even polynomials are extended to whole Gaussian moments; their finite
tails are at most $C_J e^{-t\mathcal L_+^2/400}$.

The remainder proof is incomplete until the rest of the canonical
integral is paid. In \eqref{eq:positive-heat}, split at real parts
$2$ and $-1$, reflect the left arm using $\xi(w)=\xi(1-w)$, and
keep the original bounded middle. Shift the two right rays to their
saddles with their oriented connectors. Each shifted ray and connector
stays in a half-plane where the full ordinary Dirichlet series
converges. Both actual infinite ends, the noncentral parts, connectors,
middle, and all $n\ge2$ are retained in the error account. A gamma
integrand is never extended through its poles to an invented whole-line
contour. These payments give the last two terms of
\eqref{eq:positive-budget}. A sufficient eventual absorption includes
$\kappa_a t\ge2J+4$ and an order-dependent logarithmic onset.
This is why the full bound must precede its elegant simplified form.

\subsection{The residual recurrence and the direct time estimate}
The approximation has a nonzero heat residual, with a simple exact
description. In the saddle coordinate $p$ define
\begin{equation}
 \mathcal Q_+=\frac1{4B^2}\partial_{pp}
 +\frac{t\alpha''}{4B^3}\partial_p
 +\frac{t\alpha'''}{16B^3}
 +\frac{5t^2(\alpha'')^2}{64B^4}.
 \label{eq:positive-operator}
\end{equation}
Exact differentiation of \eqref{eq:positive-saddle} gives
\begin{equation}
 \Sc^{-1}(\partial_t-\tfrac14\partial_{ww})(\Sc F)
             =(\partial_t|_p-\mathcal Q_+)F,
 \qquad
 \partial_t|_p C_{2j+2}^+=\mathcal Q_+C_{2j}^+.
 \label{eq:positive-transport}
\end{equation}
One proves the coefficient identity with finite Gaussian moments and
$\rho=t/B$: differentiation uses $\rho_t|_p=B^{-2}$ and
$\rho_p=\alpha''\rho^2/2$. Thus
\begin{equation}
 \Res P_{2J}=-\sum_{w=s,1-s}\Sc_t(w)\mathcal Q_+C_{2J}^+,
 \qquad
 (H-P_{2J})_t=-(H-P_{2J})_{zz}-\Res P_{2J}.
 \label{eq:residual-telescope}
\end{equation}
The coefficient and transport proofs are in the ansatz companion's
\src{ansatz/manuscript/07_positive_hierarchy.tex}{\texttt{prop:positive-coefficients}, \texttt{prop:positive-residual}}.
The final grade is not identically zero; in particular it is eventually
nonzero at $z=0$. Formula~\eqref{eq:residual-telescope} supplies the
time estimate without differentiating a big-$O$ bound.

\subsection{The improvement has two useful faces}
For $X\gg t\mathcal L_+$, the allowance is
$b_J\asymp_J t^{J+2}X^{-2J-3}$. In the time-dominated or comparable
region it is $O_J(t^{-J-1}\mathcal L_+^{-2J-3})$. The next allowance satisfies
\begin{equation}
 \frac{b_{J+1}}{b_J}=O_J(t/\Ns^2)
              =O_J(1/(t\log^2(2+t))).
 \label{eq:order-gain}
\end{equation}
This compares proved positive bounds, not successive signed errors.
At each fixed height, a single arm's coefficient has asymptotic form
\[
 C_{2j}^+\sim\frac{2^{2j-1}(2j)!}{(j+1)!}
                     t^{-j}\log^{-2j-1}t\quad(j\ge1).
\]
The first omitted correction explains the field scales
$t^{-1}\log^{-3}t$, $t^{-2}\log^{-5}t$, and $t^{-3}\log^{-7}t$
for $P_0,P_2,P_4$. These sharpen the field description; existing
coarser physical-transfer theorems do not improve automatically.

\subsection{A compact-positive refinement reaches the switch}
Fix $I=[a,T]\subset[32,\infty)$, with finite $T$, and let $X\to\infty$.
The actual saddle equation gives
\[
 p-w=O_T(\log X),\quad |p|=X/2+O_T(\log X),\quad
 \Re(p-w)=\frac t4\log(X/(4\pi))+O_T(\log X/X).
\]
Using this position instead of the coarser all-height gate improves
the complete arithmetic allowance to $C_I X^{-t\log2/4}$.

\begin{theorem}[Compact positive field bridge]
\label{thm:compact-positive}
For each fixed $J,I$ and fixed buffers,
\begin{align}
 |\partial_z^d(H-P_{2J})|&\le C_{J,I,d}\A(1+\log X)^d
                      [X^{-2J-3}+X^{-t\log2/4}],&&0\le d\le3,\\
 |(H-P_{2J})_t|&\le C_{J,I}\A(1+\log X)^2
                      [X^{-2J-3}+X^{-t\log2/4}],\\
 |\Res P_{2J}|&\le C_{J,I}\A X^{-2J-3}.
 \label{eq:compact-positive}
\end{align}
In particular $P_2$ has pure $X^{-5}$ field allowance on every fixed
$[32,T]$, and $P_4$ has pure $X^{-7}$ allowance on $[63,65]$, with
the displayed jet and time losses.
\end{theorem}
The numerical inequalities are $8\log2>5$ and $63\log2/4>7$.
This is a new compact-domain proof, not substitution of $64$ into an
uninstantiated eventual onset. At higher orders the arithmetic term
must remain whenever its exponent does not pay the local precision.
It is an upper-bound budget, not a proved lower bound or a sharpness
claim. The full compact proof, finite-endpoint forcing and optional
splice are in the ansatz companion's
\src{ansatz/manuscript/07_positive_hierarchy.tex}{\texttt{thm:compact-positive}},
\src{ansatz/manuscript/07_positive_hierarchy.tex}{\texttt{prop:elementary-compact}}
and \src{ansatz/manuscript/07_positive_hierarchy.tex}{\texttt{prop:positive-overlap-splice}},
with original source Note 20 \cite{bridge}.

For comparison, an elementary corrected carrier is
\begin{equation}
 F_1=\sum_\pm C_\pm(1+D_\pm),\quad
 D_\pm=\frac{\pm\ii}{z}
  \left\{\frac1{24}-\frac t4-\frac{7t}4b_\pm-\frac{t^2}8b_\pm^2\right\},
 \quad b_\pm=L_z/2\pm\ii\pi/4.
 \label{eq:elementary-correction}
\end{equation}
With $\Ls=1+\log X$, the new complete comparison on $[3/4,33]$
is
\begin{align}
 B_1(t,X)&=\Ls^4/X^2+X^{-t\log2/4},\notag\\
 |\partial_z^j(H-F_1)|&\le C_j\Eo\Ls^jB_1,&&j\le3,\notag\\
 |(H-F_1)_t|&\le C_t\Eo\Ls^2B_1,\label{eq:intermediate-band}\\
 |\Res F_1|&\le C\Eo\Ls^4/X^2,
 &|\partial_z\Res F_1|&\le C\Eo\Ls^5/X^2.\notag
\end{align}
Its proof uses the complete fixed-endpoint canonical split and a
corrected local gamma calculation, retaining all arithmetic and the
actual middle and infinite ends \cite{intermediate}. It is integrated
as the ansatz companion's
\src{ansatz/manuscript/08_positive_bridges.tex}{\texttt{thm:intermediate-positive}}.
On $[3/4,1]$, comparison of $F_0$ and $F_1$ through $H$ pays the sum
of their complete allowances; comparison of $F_N$ and $F_1$ pays
$\epsilon_X+B_1$ with $\Ls^j$ for jets and $\Ls^2$ for time.
The omitted-arithmetic term in $B_1$ is retained. On $[32,33]$ this
is precisely the existing elementary $F_1$ formula; the common buffered
jet/time premises are checked, and $\A\asymp\Eo$. That identifies
its overlap with the compact fixed-order hierarchy without a new splice.
The truncated gamma arm's finite-endpoint forcing is separately paid;
it is not treated as an exact heat solution.

For a time profile $\chi$, any optional splice obeys
\begin{equation}
 \Res[\chi P_{2J}+(1-\chi)F_1]
 =\chi\Res P_{2J}+(1-\chi)\Res F_1+\chi'(P_{2J}-F_1).
 \label{eq:positive-splice}
\end{equation}
It inherits the coarser carrier error unless an additional cancellation
is proved. Choosing fixed $T>\max(65,t_{*,J})$ gives a high-height
overlap with the all-height theorem. The remaining bounded-height
compact core is still a separate approximation task.

\section{The physical observables: exact equations before approximation}
\label{sec:physical}

\subsection{One prescribed regulator and one fixed spatial filter}
Let
\begin{equation}
 k(v)=\frac{315}{256}(1-v^2)^4\mathbf1_{|v|\le1},\quad
 k_\ell(x)=\ell^{-1}k(x/\ell),\quad K(x,y)=k_\ell(x)k_h(y).
 \label{eq:kernel}
\end{equation}
The widths are fixed; there is no time average. The kernel is $C^3$
with integrable weak derivatives through order four. The global theorem
below specializes to $\ell=h=1$.

For positive time let $\ell_s$ be the large real source saddle and set
\begin{align}
 4\pi e^{4\ell_s}&=2t\ell_s+9,
 &\kappa_s&=8t\ell_s-2t+36,\\
 M_0&=2\pi^2e^{t\ell_s^2+9\ell_s-\pi e^{4\ell_s}}
                  \sqrt{2\pi/\kappa_s},
 &\epsilon_+(t,x)&=M_0\kappa_s^{-1}e^{-x^2/(2\kappa_s)}.
 \label{eq:source-saddle}
\end{align}
With $C$ from \eqref{eq:negative-gaussian}, fix $\eta>0$ and
$\epsilon_b(t)=\eta C(1+t^2)^{-1/4}$. Let
\[
 b_*(v)=\frac{e^{-1/v}}{e^{-1/v}+e^{-1/(1-v)}}\quad(0<v<1),
\]
extended by zero and one at the respective ends, and put
$\chi(t)=b_*((t-64)/64)$. The physical regulator is
\begin{equation}
 \epsilon_g(t,x)=
 \begin{cases}\epsilon_b(t),&t\le64,\\
 \epsilon_b(t)^{1-\chi(t)}\epsilon_+(t,x)^{\chi(t)},&t>64,
 \end{cases}
 \qquad q=\epsilon_g^2,\quad r=\log q.
 \label{eq:physical-regulator}
\end{equation}
The numbers $64$ and $128$ are prescribed switch times. The analytical
floor $32$ has a different origin. None is claimed to be an intrinsic
critical time of $H$.

For a holomorphic field $F$ define
\begin{align}
 D&=|F|^2+q,& S&=\log D,& w&=q/D,\\
 \rho&=\frac{q|F_z|^2}{\pi D^2},&
 j&=\frac{qF_{zz}\overline{F_z}}{\pi D^2},&
 U&=K*S,\quad W=K*w,\\
 R&=K*\rho,&J&=K*j,&
 Q&=K*(S_x^2-S_y^2)-(U_x^2-U_y^2).
 \label{eq:physical-definitions}
\end{align}
The complex current represents the vector $(\Re J,\Im J)$.
$Q$ is the actual difference of two gradient variances. It is not a
chosen closure term or a function of density variance alone.

\subsection{Adaptive potential, curvature and current}
Write $\Dt=\partial_x^2-\partial_y^2$ and define
\begin{align}
 \mathcal B_\epsilon&=w\{r_t+\tfrac12(r_{xx}-r_{yy}+r_x^2-r_y^2)\},\\
 4\pi C_{\rm reg}&=2\operatorname{div}(w\nabla r)-w\Delta r
                                     -w(1-w)|\nabla r|^2,\\
 \mathcal S_\epsilon&=(1-2w)(r_t\rho+\nabla r\cdot\mathbf j).
 \label{eq:regulator-corrections}
\end{align}
For the exact heat field $H$, differentiation and convolution give
\begin{align}
 U_t&=-\tfrac12(\Dt U+U_x^2-U_y^2+Q)+K*\mathcal B_\epsilon,
 &\Delta U&=4\pi(R+K*C_{\rm reg}),\label{eq:physical-potential}\\
 R_t&=-\tfrac12\Dt R-\tfrac1{8\pi}\Delta(U_x^2-U_y^2+Q)
       +\tfrac1{4\pi}\Delta(K*\mathcal B_\epsilon)\notag\\
       &\hspace{1em}-\partial_t(K*C_{\rm reg})-\tfrac12\Dt(K*C_{\rm reg}),
 \label{eq:physical-density}\\
 R_t+\operatorname{div}J&=K*\mathcal S_\epsilon.
 \label{eq:physical-current}
\end{align}
For example, the first curvature identity follows by differentiating
$\log(|H|^2+q)$, using holomorphy, and collecting the $q$ derivatives.
Taking the potential equation's Laplacian proves
\eqref{eq:physical-density}; this retains both the time and signed
spatial derivatives of curvature. The full algebra is in \cite{parent}.

For constant $q$ these regulator terms vanish. In that special case,
\[
 \Delta Q=2\{\operatorname{Var}_K(S_{xx})-\operatorname{Var}_K(S_{yy})\}
 +8\pi\{\Cov_K(S_x,\rho_x)-\Cov_K(S_y,\rho_y)\}.
\]
A small-width, even-filter expansion gives a local gradient approximation
to $Q$ with fourth-order width error under a bounded $C^7$ premise.
That consistency result is not a stability theorem for a closed PDE.
Changing the filter widths with height would add commutators.

\subsection{Weak primitives preserve cancellation near zeros}
For a compactly supported real $C^2$ test $\psi$, the curvature identity
can be written
\begin{equation}
 4\pi\langle\rho,\psi\rangle
 =\int S\Delta\psi+2\int w\nabla r\cdot\nabla\psi
       +\int\{w\Delta r+w(1-w)|\nabla r|^2\}\psi.
 \label{eq:weak-density}
\end{equation}
This replaces a sharply peaked positive density by logarithmic and
bounded primitives. Near a simple zero, unit factorizations and
integrable poles then control signed current and stress expressions.
The mean product $U_x^2-U_y^2$ must still be paid separately with
bounds on the actual means. Separately estimating large positive pole
squares can destroy the cancellation that makes the estimate useful.

For a holomorphic approximate field let $E=F_t+F_{zz}$ and
$\Xi_E=2\Re(\overline F E)/D$. The potential equation acquires
$K*\Xi_E$ and the density equation acquires
$\Delta(K*\Xi_E)/(4\pi)$. The current balance acquires
\begin{equation}
 K*D_E,\qquad
 D_E=\frac{2q\Re(E_z\overline{F_z})}{\pi D^2}
             -4\rho\Re(E\overline F/D).
 \label{eq:holomorphic-defect}
\end{equation}
These exact identities identify the transfer obligations. They do not
by themselves prove that a small envelope error makes every physical
quantity small: direct inverse-$q$ estimates can be extremely costly.

\subsection{An explicit clock, with differentiated transport}
The complete source probability law, proportional to $e^{tu^2}\Phi(u)$,
has mean $\mu$ and variance $v$. Its score identity, with the half-line
endpoint retained, yields
\begin{equation}
 |\mu-\ell_s|=O(\kappa_s^{-1}),\quad
 |v-\kappa_s^{-1}|=O(\kappa_s^{-3/2}),\quad
 |\mu'/\mu-2/\kappa_s|=O(\kappa_s^{-3/2}).
 \label{eq:score-refinement}
\end{equation}
Since $\ell_s'=2\ell_s/\kappa_s$, the locally finite sums
\begin{align}
 d_j&=(j+\tfrac12)\pi/\ell_s,\qquad
 R_L(t,x,y)=k_h(y)\sum_{j\in\mathbb Z}k_\ell(x-d_j),\\
 J_L(t,x,y)&=-\frac2{\kappa_s}k_h(y)
                       \sum_{j\in\mathbb Z}d_j k_\ell(x-d_j)
 \label{eq:clock}
\end{align}
satisfy $\partial_tR_L+\partial_xJ_L=0$ exactly. Polynomial
periodization gives equivalent finite Bernoulli-polynomial formulas;
there is no omitted Fourier-phase tail. On the proved inner window of
height scale $\kappa_s^{1/3}$, with the stated buffers, the density
error is $O(\kappa_s^{-1/2})$ and the current and direct time errors
are $O(\ell_s\kappa_s^{-7/6})=o(\ell_s/\kappa_s)$.
The clock has not been assigned a potential or covariance reproducing
the entire physical tuple. Its overlap with the full saddle retains
the separate existing transfer allowance \cite{parent}.

The geometry of a block of $m$ clock roots is also explicit:
\begin{equation}
 \Delta_m^{\rm clock}
   =\left(\frac\pi{\ell_s}\right)^{m(m-1)}
                         \prod_{k=1}^{m-1}(k!)^2.
 \label{eq:clock-vandermonde}
\end{equation}
If $\mu a_i/\pi=c+i+e_i$, $|e_i|\le e<1/2$, and
$|\mu-\ell_s|\le e_s<\ell_s$, pairwise comparison gives
\begin{equation}
 \left|\log\frac{\Delta_m}{\Delta_m^{\rm clock}}\right|
 \le\frac{m(m-1)e_s}{\ell_s-e_s}
   +\frac{4e}{1-2e}\{mH_{m-1}-(m-1)\}.
 \label{eq:vandermonde-error}
\end{equation}
For the proved growing block, $m=O(\ell_s\kappa_s^{1/3})$,
$e=O(\kappa_s^{-1/2})$ and $e_s=O(\kappa_s^{-1})$, this tends to zero.
The pair Laplacian with weights $(a_i-a_j)^{-2}$ has corresponding
Loewner bounds with factors $(\mu/\ell_s)^2(1\pm2e)^{-2}$,
including its constant zero mode. These are geometric comparisons,
not a closed zero-motion law. A factorization $H=U_0\prod_i(z-a_i)$
retains the unit drift $2(U_{0,z}/U_0)(a_i)$ in the root equation;
unfolding adds the derivative of the unfolding scale.

\subsection{The regulator layer and the importance of a named center}
The earlier outer analysis uses the one-arm linear carrier
\[
 A_{\rm lin}(t,X)=|M(s_X)e^{t\alpha(s_X)^2/4}|,
 \qquad b_{\rm lin}=\log(\epsilon_g/A_{\rm lin}),\quad
 \omega\sim\tfrac14\log(X/(4\pi)).
\]
Its layer location is $d=(b_{\rm lin})_+/\omega$, and its leading
density is
\begin{equation}
 R_{\rm layer}=\frac{\omega}{2\pi}
                       \{k_h(y-d)+k_h(y+d)\},
 \label{eq:layer}
\end{equation}
with the stated $O(1)$ absolute error in the source theorem. Define
the exact center by $b_{\rm lin}(t_c(X),X)=0$. Then
\begin{align}
 t_c-64&=\frac{64}{\log X+c_0+O(1/\log X)},
 &c_0&=\log\frac{4e}{\pi\kappa_s(64)},\\
 t_X(\sigma)&=t_c(X)-\frac{128\sigma}{\pi X\log X},\\
 \frac{R_H(t_X(\sigma),X,y)}{\log X}
 &=\frac{k_h(y-\sigma_+)+k_h(y+\sigma_+)}{8\pi}
                           +O(1/\log X).
 \label{eq:resolved-layer}
\end{align}
The scaled-time statement is uniform for bounded $\sigma$ on its
declared spatial domain. At the exact center,
\begin{equation}
 \partial_t b_{\rm lin}(t_c,X)
   =-\frac\pi{512}X(\log X)^2[1+O(1/\log X)].
 \label{eq:center-sensitivity}
\end{equation}
The center shift $O(1/\log X)$ is far larger than the rounding width
$O(1/(X\log X))$. Substituting its coarse asymptotic expansion into
the resolved law is therefore invalid. The full-saddle and corrected centers retain their separate names.
Put $t_0=t_c$, $\lambda=\log X$, $w_X=128/(\pi X\lambda)$,
and keep $B_0=\partial_t b_{\rm lin}(t_0,X)$ exact. For a named one-arm
modulus $A$, set $\Delta_A=\log(A/A_{\rm lin})$ and define $t_A$ by
$\log(\epsilon_g/A)=0$. The paid comparison yields \cite{centers}
\begin{align}
 t_{A_0}-t_0&=-\frac{7\pi t_0}{16XB_0}
                  +O((X^3\lambda)^{-1}+X^{-4}),\notag\\
 t_A-t_0&=\frac{\pi t_0^2L}{32XB_0}+O(\Ls^2/X^3),
                    && A=A_1,A_J,A_S,A_4,A_{\exp},\label{eq:named-centres}\\
 t_{2A}-t_A&=\frac{\log2}{B_0}+O((X^2\lambda^2)^{-1}).\notag
\end{align}
Here $A_0=|C_+|=A_o/2$, $A_1=|C_+(1+D_+)|$,
$A_J$ is the fixed truncated-gamma arm, $A_S=|\Sc_t(s_X)|$,
$A_4=|\Sc_t(s_X)(1+C_2^++C_4^+)|$, and $A_{\exp}$ is the
separately defined exponential rewrite of the corrected modulus.
The elementary center is later by order $1/(X^2\lambda^2)$; corrected
centers are earlier by order $1/(X^2\lambda)$. Even the distinct
one-arm/cosine factor two moves the center by only
$O(1/(X\lambda^2))$. All are $o(w_X)$. Exact flat-switch derivatives
pay this comparison uniformly on a fixed multiple of the rounding
neighborhood; coarse differentiation of a center-location allowance
would not. The detailed argument is the ansatz companion's
\src{ansatz/manuscript/09_switch_consumer.tex}{\texttt{prop:named-switch-centres}}.
The older simplified-layer law retains $O(1)$ absolute and
$O(1/\log X)$ normalized accuracy.

\subsection{A complete specialized fixed-\texorpdfstring{$P_4$}{P4} physical consumer}
\begin{theorem}[Compact switch tuple]
\label{thm:p4-physical}
Fix $\ell,h,\eta>0$, $K=k_\ell k_h$ and exactly
$q=\epsilon_g^2$. On $63\le t\le65$, $X=|x|\ge X_0$ and
$|y|\le1$, let $P=P_4$ and $e=X^{-7}$. For the full physical
observables in \eqref{eq:physical-definitions} and
\eqref{eq:regulator-corrections}, write $\delta$ for $H-P$ at the
observable level, and $\mathcal C=K*C_{\rm reg}$. Finite
$C,X_0$ depending on the fixed data and named source constants give:
\begin{center}\small
\begin{tabular}{p{.52\textwidth}cc}
\toprule Outputs & Spatial orders & Absolute allowance\\\midrule
$\delta U,\delta W$ & $\le4$ & $Ce$\\
$\delta R,\delta\mathcal C$ & $\le2$ & $Ce$\\
$\delta J,\delta Q$, including $\Delta\delta Q$ & $\le2$ & $Ce\Ls^2$\\
$\delta J$ used internally for source balance & $\le3$ & $Ce\Ls^2$\\
$\delta U_t,\delta W_t,\delta R_t,\delta\mathcal C_t$ & $\le2$ & $CeX\Ls^2$\\
$K*(b_{\rm reg}\delta w),\ K*\delta\mathcal S_\epsilon$ & $\le2$ & $CeX\Ls^2$\\
$K*\Xi_E$ (heat-potential correction) & $\le4$ & $Ce$\\
$K*D_E$ (density/current correction) & $\le2$ & $Ce$\\\bottomrule
\end{tabular}
\end{center}
Here $b_{\rm reg}=r_t+(r_{xx}+r_x^2)/2$,
$E=P_t+P_{zz}$, and $\Xi_E,D_E$ are the actual residual corrections
of Section~\ref{sec:physical}. All means, weak masses and actual signed
covariance are retained. $P$ is holomorphic and unjoined; compatibility
and field-join defects vanish, while $E$ is generally nonzero.
\end{theorem}
The complete specialized proof is in the ansatz companion's
\src{ansatz/manuscript/09_switch_consumer.tex}{\texttt{thm:fixed-p4-consumer}},
from Note 23 and Assessment 17 \cite{p4-physical}. It consumes the
parent's exact physical identities and weak primitive lemmas. Corrected
phase disks acquire every $H/P$ collar and zero-free analytic unit anew.
Subtracting weak logarithms, bounded weights and signed poles before
integration pays $\|\delta S\|_1+\|\delta w\|_1=O(e)$,
$\|\delta V_m\|_1=O(e\Ls^2)$ and residual
$\|Z_m\|_1=O(e)$. The actual full mean products are bounded separately,
then combined with the signed stress to obtain $Q$ and $\Delta Q$.

The exact flat switch is written $r=a(t)-\beta(t)x^2$.
Where $\beta X\le B_*$ its true time derivatives cost
$O(X\Ls^2)$; where $\beta X>B_*$ the regulator-weighted terms are
exponentially small. The first inequality is an estimate regime;
$B_*/X\le1/2$ is an eventual onset gate. Direct time identities,
rather than derivatives of allowances, prove $R_t$ and curvature-time
rows; current through order three then pays the actual signed source.
All tests use weak kernel derivatives through four. The source constants
and finite height onset remain symbolic. This high-height compact
consumer supplies no new bounded-core, small-positive physical or global
joining theorem, and does not sharpen the simplified-layer rate.

For each fixed time, the existing physical tail law is
$R_H/\log X\to0$ for $t\le64$ and $R_H/\log X\to k_h(y)/(4\pi)$
for $t>64$ on $|y|<h$. This is a consequence of the prescribed
regulator. It does not identify a critical time of the unregulated field.

\section{A finite construction for all time and all height}
\label{sec:global}

\subsection{Spatial joining changes the equations}
Let $G_s$ be the entire exact-heat field of a fixed nonnegative finite
source, and let $P=P_0$ be the established leading full saddle.
In a paid model gap join them by
\begin{equation}
 F=G_s+\lambda(t,x)(P-G_s),\qquad D_0=P-G_s.
 \label{eq:spatial-join}
\end{equation}
The real profile makes $F$ nonholomorphic. With $B_c=F_{\bar z}$ and
$E=F_t+F_{zz}$, its exact defects are
\begin{align}
 B_c&=\tfrac12\lambda_xD_0,&
 E&=\lambda\Res P+\lambda_tD_0+\lambda_x(D_0)_z
                                          +\tfrac14\lambda_{xx}D_0,\\
 (B_c)_z&=\tfrac14\lambda_{xx}D_0+\tfrac12\lambda_x(D_0)_z,&
 (B_c)_{\bar z}&=\tfrac14\lambda_{xx}D_0,\\
 (B_c)_{zz}&=\tfrac18\lambda_{xxx}D_0+\tfrac12\lambda_{xx}(D_0)_z
                                          +\tfrac12\lambda_x(D_0)_{zz}.
 \label{eq:spatial-defects}
\end{align}
Time derivatives of these terms are retained when the gap moves.
$P$ is evaluated only on its established exterior branch until its
all-height late-time theorem applies.

Use the same algebraic observables \eqref{eq:physical-definitions}
with $a=F_z$, $c=F_{zz}$ and set
$V_m=w^mc\overline F/D$, $Z_m=w^mE\overline F/D$. Then
\begin{align}
 T_C&=\frac{q|B_c|^2}{\pi D^2}
  +\frac2\pi\Re\left(\frac{(B_c)_z\overline F}{D}
                       -\frac{aB_c\overline F^2}{D^2}\right),\\
 C_{\rm tot}&=C_{\rm reg}+T_C,\qquad
 \Xi=8\Re\frac{a\overline{B_c}}D
                  +4\Re\frac{F\overline{(B_c)_{\bar z}}}D,\\
 T_J&=-\frac{(B_c)_{zz}\overline F}{D}
                    +\frac{cB_c\overline F^2}{D^2},
 &T_U&=\Xi/2+2\Re Z_0.
 \label{eq:nonholomorphic-corrections}
\end{align}
For the physical $q=q(t,x)$, direct differentiation gives
\begin{align}
 \Delta S&=4\pi(\rho+C_{\rm tot}),
 &\pi j&=\partial_{\bar z}V_0+(\partial_{\bar z}r)V_1+T_J,\\
 S_x^2-S_y^2&=4\Re V_0-\Dt S+(r_{xx}+r_x^2)w+\Xi,
 &S_t&=-2\Re V_0+r_tw+2\Re Z_0.
 \label{eq:nonholomorphic-primitives}
\end{align}
The potential equation acquires $K*T_U$ and the density equation
acquires $\Delta(K*T_U)/(4\pi)$, with $C_{\rm tot}$ replacing
$C_{\rm reg}$. The current equation is
\begin{align}
 R_t+\operatorname{div}J&=K*\mathcal S_\epsilon+K*(D_E+D_B),\\
 D_E&=\frac{2q\Re(E_z\overline a)}{\pi D^2}-4\rho\Re Z_0,\\
 D_B&=\frac{2q\Re(c\overline{(B_c)_z})}{\pi D^2}
              -\frac{4q\Re(F\overline{B_c}c\overline a)}{\pi D^3}.
 \label{eq:nonholomorphic-current}
\end{align}
These identities follow from the actual derivatives of $D$; no
holomorphic premise is imposed on the joined field. Its algebraic
density is not asserted to count zeros in the joining bands.

\subsection{The complete approximation contract}
Write $\mathcal C=K*C_{\rm tot}$,
$\overline{\mathcal B}=K*\mathcal B_\epsilon$ and
$\overline{\mathcal S}=K*\mathcal S_\epsilon$. For a real spatial
multi-index, the established scalar contract is
\begin{equation}
 \begin{array}{c|c}
 \text{quantity compared with its value for }H&\text{spatial orders}\\\hline
 U,W&0\text{--}4\\
 R,\mathcal C,J,Q,\overline{\mathcal B}&0\text{--}2\\
 U_t,W_t,R_t,\mathcal C_t,\overline{\mathcal S}&0\text{--}2\\
 K*T_U\quad(\text{zero for }H)&0\text{--}4
 \end{array}
 \label{eq:global-tuple}
\end{equation}
Real and imaginary current components count separately. Internal current
order three pays source order two. The contract includes $\Delta Q$
and both curvature derivatives in the density equation. A scalar
tolerance applies to each component; an equation containing a sum
inherits the sum of its component allowances. No second time derivative
or moving observation boundary is included.

\begin{theorem}[Existing finite all-time construction]
\label{thm:global}
Fix $\eta>0$, $K=k_1k_1$ and $0<\theta\le1$. There are explicitly
specified finite parameters and one nonnegative finite source for which
the smooth field below approximates every scalar of
\eqref{eq:global-tuple} within $\theta/2$, uniformly for
$t,x\in\R$ and $|y|\le1$:
\begin{equation}
 F(t,z)=\begin{cases}
 N_{-t}(z),&t\le-2T_-,\\
 (1-v_-)N_{-t}+v_-G_s,&-2T_-\le t\le-T_-,\\
 G_s,&-T_-\le t\le63,\\
 G_s+v_0\lambda(P-G_s),&63\le t\le64,\\
 G_s+\lambda(P-G_s),&64\le t\le T_+,\\
 G_s+[v_++(1-v_+)\lambda](P-G_s),&T_+\le t\le T_++1,\\
 P,&t\ge T_++1.
 \end{cases}
 \label{eq:global-field}
\end{equation}
Here $v_-=b_*((t+2T_-)/T_-)$, $v_0=b_*(t-63)$, and
$v_+=b_*(t-T_+)$. The physical regulator and filter are unchanged.
\end{theorem}
This is the parent companion's
\src{manuscript/23_all_time_assembly.tex}{\texttt{thm:all-time}},
integrating Notes 54--55; it is not a new consequence of the corrected
positive hierarchy \cite{global,conditioning}.

\begin{proof}[Structure of the finite construction]
Choose the late-time physical coefficient table and $T_+$ first.
Then fix $b=T_++1$ and build a compact exterior table on $[63,b]$.
Its constants retain their $b$ dependence. Select $T_-$, a sufficiently
large exterior phase gap, a frozen scaled width and the moving gap
trajectory. Finally choose one source satisfying the minimum of the
finitely many positive accuracy demands. This order avoids using
compact constants to select the late-time constants that define them.

On the central region a positive regulator floor pays direct source
transfer, including multiple zeros. Outside it, direct $G_s/H$ collars,
unit factorizations and integrable poles replace a growing inverse-floor
bound. The spatial transition lies in a model gap with a common field
floor. Complete convolution supports are divided for estimates only;
their moving boundaries are not differentiated as new fields.
Whole-support positive primitive allowances are added before applying
the signed identities and actual covariance mean product.

The compact exterior allowance is
$O_b(\log^8X_0/X_0)$, with a sharper existing
$O(\log^6X_0/X_0^2)$ table on $[95,97]$. Source refinement is
simultaneous with the choice of $X_0$. Flat temporal profiles pay the
two endpoint attachments and the initial exterior activation. The
additional $v_+'(P-G_s)$ and moving compatibility terms are included.
The finite positive sums fit the scalar reserve $\theta/2$.
At any time only the active row is compared, so no multiplier by the
number of rows is needed. Detailed finite coefficients and gates are
in the parent proof \cite{global}.
\end{proof}

At negative infinite time the density, current, density-time/covariance,
and $\Delta Q$ errors have the established orders
$T^{-2}$, $T^{-5/2}$, $T^{-3}$, and $T^{-4}$ respectively, with the
adaptive regulator. At positive infinite time the existing static
$U,W$ errors are $O(t^{-1/2}\log^{-2}(2+t))$ and the retained
evolution and covariance errors are $O(t^{-1/2})$, uniformly in
height. The leading saddle's heat defect has its separate $t/\Ns^3$
scale. Errors tend to zero; the physical density itself need not have
a finite limit.

Replacing $P_0$ by $P_{2J}$ in \eqref{eq:global-field} is a promising
refinement, but requires new collars, means, physical-transfer and
joining estimates. The global theorem above retains its established
model and rates until the consumers for that new assembly are proved. The specialized $P_4$ theorem on $[63,65]$ pays its unjoined compact tuple; neither it nor the positive field atlas has been inserted into a new global spatial/time join. The wider collars, moving gaps and all new compatibility terms would require their own complete budget.

\subsection{Why a finite source is a bridge, not a universal endpoint}
Three scoped obstructions clarify the architecture. For a fixed
constant regulator and fixed finite time, the canonical density tends
to zero at infinite height on bounded transverse strips; it cannot
uniformly realize a clock density growing like $\log X$. A fixed
compact source supported in $[0,U]$ has positive-time growth controlled
by $e^{tU^2}$, whereas the true source saddle escapes to infinity.
Under the prescribed positive regulator its density tends to zero
while the canonical coarse density has scale $\ell_s k_h/\pi$.
This excludes a fixed compact source as the positive endpoint model,
not finite sources on bounded upper-time families.

A further potential obstruction is proved for the specified class of
continuous piecewise-linear compact sources with positive terminal
value, at fixed $t>64$ and the stated off-axis locations. Its algebraic
endpoint tail mismatches the canonical exponential tail and Gaussian
regulator. It is not an exclusion of every finite representation.
The full-source absolute derivative hierarchy remains a fallback on
compact domains, but does not supply an elegant normalized global
asymptotic description. These distinctions explain the three-mechanism
construction rather than undermine it.

\section{Positive source compression and finite evaluation}
\label{sec:effectivity}

\subsection{Approximate a square root, then square}
The fixed-source ingredient can be made substantially smaller without
changing the global physical theorem. On
\[
 \mathcal D=\{u\in\C:\Re u>-1/16,\ |\Im u|<1/16\}
\]
the complete canonical series admits the factorization
\begin{align*}
 \Phi(u)&=A(u)(1+\mathcal R(u)),\\
 A(u)&=\pi e^{5u}(2\pi e^{4u}-3)e^{-\pi e^{4u}},\qquad
 |\mathcal R(u)|<1/2.
\end{align*}
Both factors used for their logarithms lie in a right half-plane.
Hence
\begin{equation}
 h(u)=\exp\{[\log\pi+5u+\Log(2\pi e^{4u}-3)
                           -\pi e^{4u}+\Log(1+\mathcal R(u))]/2\}
 \label{eq:source-root}
\end{equation}
is an explicitly selected analytic square root, positive on the real
half-line, with $|h|<9$ on $\mathcal D$. The tail bound includes
every arithmetic mode; real positivity alone would not have supplied
this complex branch.

Approximate $h$ on $[0,U]$ by a polynomial with real dyadic coefficients
in a normalized Chebyshev basis, round those coefficients with a paid
error, and set $f=p^2$ on $[0,U]$, zero beyond it. Positivity is then
exact. The full tail above $U$, interpolation error, squaring error
and coefficient error give
\begin{theorem}[Compressed positive source]
\label{thm:source-compression}
Fix a finite upper time $b$. For $0<\varepsilon\le1$ there is a
finite nonnegative polynomial source with exact-heat field $G_s$ such that
\begin{equation}
 |\partial_z^j(G_s-H)|\le51\varepsilon/256,
 \qquad0\le j\le3,\quad t\le b,\quad x\in\R,\quad |y|\le4.
 \label{eq:source-four-jets}
\end{equation}
As $q_s=\log(1/\varepsilon)\to\infty$ at fixed $b$, its sufficient
degree is $O_b(q_s\log q_s)$, coefficient fractional precision is
$O_b(q_s)$, and compressed descriptor length is
$O_b(q_s^2\log q_s)$ bits. Expansion into normalized powers costs
$O_b(q_s^2\log^2q_s)$ bits. A finite coarse construction recipe has
polynomial bit cost, bounded by $O_b(q_s^{60})$ in the stated model.
\end{theorem}
The bit-cost statement uses a dyadic tolerance and a rational upper bound
for $b$, or supplied certified recipe integers. It excludes the cost of
obtaining the upstream all-time contour and transfer constants.
This is Note 58's construction \cite{compression}, now proved in the
parent companion's \src{manuscript/25_analytic_source.tex}{\texttt{thm:ac-source}, \texttt{thm:ac-construction}}. It supplies the
same four jets consumed by the global theorem. The source changes
when the requested tolerance changes, but remains fixed during one
construction and independent of the queried $t,x,y$.

The improved fixed-central conditioning argument gives, at fixed
$b,\eta,\theta$ and source comparison radius $R_c$,
\begin{equation}
 q_s=\pi R_c/8+O_{b,\eta,\theta}(\log^2R_c).
 \label{eq:source-demand}
\end{equation}
Consequently the sufficient compressed descriptor is
$O_{b,\eta,\theta}(R_c^2\log R_c)$ bits. This improves the earlier
specified clipped-cubic sufficient exponential descriptor bound.
It is an upper bound for this construction, not a lower bound,
an affordable finite crossover, or a uniform statement when $b$ grows
with the accuracy request. Earlier positive Hermite and signed-source
constructors remain valid alternatives with their recorded scopes.

\subsection{Finite source moments and removable formulas}
If $f(u)=\sum_{k=0}^d d_k(u/U)^k$ on $[0,U]$, put
\[
 M_k(a,c)=\int_0^1v^ke^{av^2+cv}\dd v.
\]
Then
\begin{equation}
 G_s(t,z)=\frac U2\sum_{k=0}^d d_k
                  [M_k(tU^2,\ii zU)+M_k(tU^2,-\ii zU)].
 \label{eq:source-moments}
\end{equation}
Spatial jets shift $k$ and introduce powers of $\pm\ii U$.
Integration by parts gives
\[
 2aM_1=e^{a+c}-1-cM_0,\qquad
 2aM_{k+1}=e^{a+c}-kM_{k-1}-cM_k\quad(k\ge1).
\]
The entire expansion
\begin{equation}
 M_k(a,c)=\sum_{r,s\ge0}
                 \frac{a^rc^s}{r!s!(k+2r+s+1)}
 \label{eq:entire-moments}
\end{equation}
provides a route across removable denominators. An infinite series is
not itself a finite evaluator: the evaluation recipe chooses a finite
truncation, bounds its tail and pays cancellation in the coefficient sum.

\subsection{Selecting a phase index without an equality oracle}
The global construction may enlarge its onset. Note 60 uses that
freedom to choose a dyadic starting height with certified distance
from an integer phase. For the canonical phase at time $64$,
\[
 \Theta(64,x)=\frac x4(\log(x/(4\pi))-1)
              +4\pi\log(x/(4\pi))+\frac{7\pi}8+E(x),
 \qquad |E(x)|\le2^{16}(1+\log x)^2/x
\]
at the declared large onset. If $Q_j=\Theta(64,2^{K+j})/\pi$,
$j=0,1,2,3$, then
\begin{equation}
 Q_3-5Q_2+8Q_1-4Q_0=4\log2+\text{error},
 \qquad |\text{error}|<1/100.
 \label{eq:phase-stencil}
\end{equation}
The right side stays more than $19/100$ from an integer. The absolute
coefficient sum is eighteen, so at least one $Q_j$ stays more than
$1/100$ from every integer. Certified phase intervals select a candidate with distance greater than
$1/200$ from the integers and determine its exact ceiling. With $K\ge2048$ and all other
onset gates, the cost is polynomial in $K$, not in $\log K$.
This solves the phase selection needed by the construction; it is
not an algorithm for deciding equality at an arbitrary prescribed
onset. The parent companion's \src{manuscript/26_complete_evaluation.tex}{\texttt{thm:phase-selector}}
contains the finite selector proof \cite{phase}.

\subsection{The full evaluator and its cost}
Notes 59--61 treat every scalar of \eqref{eq:global-tuple}, including
the regulator, gamma saddle, profiles, source, actual covariance mean
product, continuous convolution and signed defects. One additional
local reserve can be chosen proportional to
\[
 \zeta(t,x)=\frac{\theta}{4(1+t^2)(1+x^2)},
\]
so the evaluation error does not erase the endpoint limits. It is
combined with the analytical construction's reserve component by
component, not treated as an exact arithmetic output.

The improved cost proof uses explicit pure-axis Gevrey bounds
$|\partial_i^rf|\le V A^r(r!)^3$ on each fixed query/slack box.
Rational interpolation of degree $n=O(\log(1/e))$ on
$O(An^2)$ panels per axis gives $O(A^2n^6)$ tensor nodes. Negative
interpolation weights and their amplification are paid. Shifted
Stirling with a number of shifts and terms linear in precision supplies
normalized log-gamma and fixed jets; gamma is not a unit-cost oracle.

\begin{theorem}[Polynomial additional-accuracy cost at fixed data]
\label{thm:cost}
Fix a certified descriptor and a compact rational query/slack box
$\mathcal Q$. For rational or dyadic independent scalar inputs, with
the specified canonical providers, additional absolute output precision
$p$ and query encoding length $I$ have sufficient bit bounds
\begin{equation}
 \operatorname{work}\le C_{\rm desc,\mathcal Q}(1+p+I)^{40},\qquad
 \operatorname{storage}\le C_{\rm desc,\mathcal Q}(1+p+I)^8.
 \label{eq:physical-complexity}
\end{equation}
For general computable scalar inputs let $C_{\rm ext}(P)$ and
$S_{\rm ext}(P)$ bound one request to all external scalar providers.
With $P_*=C_{\rm desc,\mathcal Q}(1+p+I)$, multiply the displayed work
bound by $1+C_{\rm ext}(P_*)$ and add $S_{\rm ext}(P_*)$ to storage.
An arbitrarily slow computable scalar's digits cannot be hidden in a
fixed constant. No uniform polynomial cost in height, time or descriptor
size is claimed.
\end{theorem}
The constants still contain inverse physical floors and can be
enormous. An application must separately acquire contour or first-jet margins, isolated zero counts/locations and initial energy, then pay buffers, complete mode counts, height onsets, separation and contour evaluation. In the moving-negative theorem even the outer sufficient quadrature count is $O(L^3)$; provider work is additional. None of these acquisitions is hidden in the fixed-data polynomial precision statement. The full evaluator and cost proofs are the parent companion's
\src{manuscript/26_complete_evaluation.tex}{\texttt{thm:complete-evaluator}, \texttt{thm:polynomial-accuracy}},
from Notes 59--61 \cite{evaluation,phase,cost}; the result
does not close the historical full-certificate or full-audit gates.

\subsection{A complete weak-filter consumer}
Notes 71--75, collected in Summary 76, extend the fixed-box theorem
above at explicitly bounded scopes \cite{weak-consumer,graded-cost,
binary-p,all-height-cost,binary-g,summary76}. Their source status is
``proved in-note from named analytical inputs; same-model internal
assessment only.'' The proof homes remain those notes; this section
records their contracts and mechanisms.

The first step is to reduce the \emph{actual} global identities to weak
primitives before making numerical holes. The inventory contains
potential and regulator means, density, current, independent direct-time
rows, signed source, actual covariance, both curvature rows, every
retained compatibility/heat-defect component and the specified traces.
The heat variation retains the nonholomorphic contribution
$-\delta_E T_C$. Covariance is formed from the signed stress and products
of the actual full means. Component covariances and positive-square
substitutes would change the observable.

On the acquired outer simple-root collars, logarithmic and simple-pole
primitives have omission costs of order $h^2$ times a logarithmic factor
and $h$, respectively. These costs, rational masks, collar guards,
normalized providers and post-filter amplification pay a complete
consumer without an inverse of the expanding outer regulator floor.
The central multiple-event region retains its fixed positive floor.
Every collar meeting the full enlarged support is included. At weak
kernel order four the spatial translation estimate has a square-root
modulus; a cube-root time modulus needs the additional time jets and
root-motion data stated in Note 71. Rational query-time evaluation does
not silently import those extra hypotheses.

Deleting disks from raw density would lose nonvanishing mass. Numerical
masks approximate the already reduced primitive integrals; they do not
redefine the physical density, filter, field or source. Note 71's
uniform-grid construction is finite but can cost $O(e^{-6})$ at error
$e$. The later graded argument supplies the accuracy-bit improvement.

\subsection{Scoped acquisition and binary-height query bounds}
Fix the existing global family whose admissible phase index is anchored
at time $64$, its canonical recipe constants, physical $q,K$, upper-time
parameter and accepted source tolerance $\varepsilon$. Choose the source
\emph{before} requesting additional evaluation error $2^{-p}$. Put
$q_s=1+\log(1/\varepsilon)$, using an upward integer bound in algorithms.
Let $I$ include rational query and relevant header encoding, and $D$ the
actual descriptor length. These parameters have different roles:
increasing $p$ does not refine the source-to-$H$ error.

\begin{theorem}[Complete costs on specified query families]
\label{thm:scoped-query-cost}
Under the canonical finite-data, buffer, margin and scalar-encoding
hypotheses of Notes 72--75, the complete retained physical tuple at its
paid orders has sufficient work/storage bounds
\[
 W\le C T_*^{80},\qquad S\le C T_*^{32},
\]
with streamed outputs and the following choices.
\begin{enumerate}[leftmargin=*]
\item On the fixed-width seam/slack family in $[95,97]$, with the fixed
central fallback, Note 72 uses $T_*=1+R+q_s+p+I$, where $R$ is a
\emph{numerical} buffered core radius. Source construction, admissible
phase selection, actual collar roots and provider cancellation are paid.
Its sufficient node bound is $CR^6(p+1+\log(2+R)+\log C)^7$.
\item On supports and full slack wholly in the guarded pure-$P$ exterior
of that slab, Note 73 uses $T_*=1+p+k+I$, where
$k=\lceil\log_2(10+\lceil|x_0|\rceil)\rceil$. Global construction
is separate; inactive source coefficients are not evaluated. The actual
residual $E_P$ is retained.
\item For \emph{every} horizontal query on $[95,97]$, Note 74 uses
\[
 A_g=1+R_g+q_s+D,\quad B_q=1+p+k+I,\quad T_*=A_g+B_q.
\]
Here $R_g$ bounds the fixed global core and its buffers. The complete
cover includes central multiple events and all band-crossing kernels.
The exterior keeps the sharper $CB_q^{80}/CB_q^{32}$ query bound.
\item On the genuine pure-$G$ plateau $[62,63]$, Note 75 uses
\[
 T_*=2+U+m+\beta+q_s+D+p+I+k,
\]
where $U,m,\beta$ are its exact normalized polynomial-source encoding
parameters. For the canonical encoding this is equivalent, up to fixed
factors, to $1+q_s+D+p+I+k$. All horizontal heights and arbitrary root
multiplicities are permitted at the fixed physical floor.
\end{enumerate}
The constants depend on the fixed geometric, normalized-margin,
physical-floor, kernel and finite-circuit recipes specified in each
source. They are not uniform as those floors vanish or time ranges
change. General external scalar providers multiply work by
$1+C_{\rm ext}(CT_*^2)$ and add $S_{\rm ext}(CT_*^2)$ to storage.
Uncached monolithic descriptor scans retain their actual cost.
\end{theorem}

\paragraph{Proof mechanisms and accounting.}
Note 72 integrates on dyadic rings: scaled area pays the simple pole
before smoothness is estimated, and large logarithmic offsets do not
enlarge positive logistic derivatives. Exact support clipping handles
the fourth kernel derivative's edge jumps. Signed quadrature weights,
source cancellation, roots, arithmetic and actual mean products are
all charged. Note 73 pays the full weak saturation replacement when
its guard holds; in the complementary regime numerical $R^2$ is bounded
by the precision demand. Reduced-phase/log-arm providers then give the
binary-height bound. Note 74 reduces primitives globally before
dispatching overlapping rational charts, so crossing a numerical chart
boundary discards no physical defect. Its dispatch radius does not
silently strengthen the already accepted source tolerance.

On $[62,63]$ the regulator is time-only and
\[
 q\ge(\eta c_{\rm base}/8)^2>0.
\]
This permits uniform fixed-floor filtering through any roots, with
$O((p+1)^6)$ nodes. The oscillatory source provider is the remaining
issue. Note 75 uses finite integration by parts at high frequency,
retaining \emph{both} source endpoints and paying the remainder; exact
endpoint-polynomial sums and equality-safe phase reduction pay
cancellation. At lower frequency the earlier entire series has only
polynomial numerical frequency. The source's zero extension is not
treated as an analytic function across its endpoint. Direct-time rows
and the nonzero regulator source remain in the tuple even though this
regional field has no heat or nonholomorphic defect.

For the all-height $[95,97]$ family the one-time construction remains
\[
 W_{\rm build}\le W_{\rm fixed}+C(1+R_g+q_s+D)^{80},\qquad
 S_{\rm build}\le S_{\rm fixed}+C(1+R_g+q_s+D)^{32}.
\]
Repeated queries add their own work and maximum streamed workspace.
Binary \emph{query} height is not a binary-global-core constructor.
The degrees $80/32$ are conservative sufficient bounds, not measured
runtime, lower bounds or evidence of practical efficiency. No evaluator
implementation or numerical certificate was executed for these notes.

\subsection{What remains after the cost update}
The complete weak consumer is no longer an unfulfilled fixed-data task.
The prepared continuation is the actual temporal blend on $[63,64]$;
on its interior $F=G+v_0\lambda(P-G)$ and its temporal defects must be
paid. It is not part of Note 75's pure-$G$ theorem. The approach through
$64<t<95$, practical constants, broader time coverage and independent
inherited-chain validation remain separate.

Endpoint applications introduce additional parameters: vanishing
observation width, shrinking regulators, multiplicities and lower-unit
conditioning. Outer simple-root puncture estimates and fixed-floor
arbitrary-multiplicity estimates do not together prove a uniform theorem
in all those parameters. At zero time the carrier is exponentially small
in height, so an absolute-source-error route to polynomially small
\emph{normalized} error requires $q_s$ of numerical-height size. This
is a limitation of that accounting route, not a lower bound for every
normalized arithmetic evaluator. The collision targets in the next
section specify a distinct consumer of the approximation machinery.

\section{Toward an economical all-region hierarchy}
\label{sec:frontier}

\subsection{A precise objective}
The eventual objective is a family of explicitly evaluable descriptions
of the same canonical field and physical observables, indexed by a
finite accuracy choice. Each member should carry its domain, positive
envelope, jets, residual, joining terms and evaluation budget. Increasing
the choice should produce a proved improvement with understandable
constants. The desired endpoint is a short formula for the error as
well as for the approximation.

This objective is a quantitative approximation programme. It is not
an RH-equivalent criterion, a new positivity implication, or a claim
to exclude collisions. Heat identities and finite recurrences apply
to much broader classes. The quantitative canonical inputs used here
are the exact theta source, its arithmetic completion, the gamma
factor and its paid sectorial estimates. Finite-product, cosine and
collision controls do not satisfy that complete canonical source
premise; their existence nevertheless prevents using the common
shape or heat equation alone as a zero-location argument.

\subsection{The next gains, in dependency order}
The six returns summarized in \src{notes/67_CHAT0025_RESEARCH_SUMMARY_AND_HANDOFF_20261026.md}{Summary 67}
and the \src{notes/68_MANUSCRIPT_UPDATE_PLAN_AFTER_RH_APPLICATION_REVIEW_20261026.md}{production plan}
are integrated at their bounded scopes. All retain ``proved in-note
from named inputs; same-model internal assessment only.'' The live
branch registers retain their wider unfinished programmes.
\begin{enumerate}[leftmargin=*]
\item \textbf{Instantiate constants and acquire lower data.}
The positive field atlas now reaches zero at high height. The fixed-$P_4$
physical switch consumer also closes. Their constants/onsets remain
symbolic. For an RH-facing consumer, first obtain an independently
proved contour or first-jet floor, identified isolated clusters and
initial energy; upper envelope errors are not lower margins.
\item \textbf{Pay a matched exterior and energy consumer.}
The conditional contour interface below supplies a better error rule,
but the complete model drift, boundary selection, reference-strip,
shrinking-index and integrated-energy payments remain obligations.
A methods pilot in an already known real-zero region is not new coverage.
\item \textbf{Choose stronger approximation from a measured need.}
The signed crossover already gives vanishing microscopic field error
at $|t|=O(L^{-2})$. Strong negative precision at $T\asymp L^{-2}$
is a separate possible task, justified by a specific consumer budget.
More late-time gamma orders or switch precision have lower priority for
this application until they improve a complete margin, onset, range or
acquisition cost. The two-arm arithmetic allowance at $t=c/L$ tends to
$2^{-c/4}$; retaining modes explains why $F_N$ is the useful near-zero
input. This is a diagnosis of that allowance, not a lower bound on error.
\item \textbf{Continue the broader physical/effectivity programme.}
Small-positive physical transfer, bounded-height compact refinement,
explicit height/time/margin costs and the complete weak-primitive filter
consumer remain separate. Expose fixed-$J$ constants before selecting
a growing order; retain the complete arithmetic population at each step.
\item \textbf{Reassemble and audit.}
Replace components of Theorem~\ref{thm:global} only after paying the
new full consumer and join contracts. Sum regional and joining budgets
without changing norms. Complete implementation/certification and an
independent inherited-chain audit are distinct milestones.
\end{enumerate}
These priorities follow the repository's
\href{run:\researchroot/result-summaries/11_newman_heat_flow_reductions_and_certificates.md}{heat-flow result index},
\href{run:\researchroot/result-summaries/12_collision_controls_and_screening.md}{collision controls}
and \href{run:\researchroot/avoid/12_term_by_term_matching_and_separate_bounds.md}{coupled-payment guidance};
the lower-margin priority does not cancel the independent all-region
approximation objectives.

\subsection{A conditional exterior-drift payment}
The application frontier is pinned to Rodgers--Tao Note 02 and its
Assessment 01 \cite{rt-interface}. The bounded RT-01 source audit is
complete; RT-03/RT-04 r2 is the separate continuation in its sole
application register. This synthesis does not execute that request or
claim later work from the concurrent application chat.

Fix one time-independent, real-symmetric contour $\Gamma$ of length
$p$ in a paid analytic chart and one fixed real interior interval $J$
at distance at least $b>0$ from it. The same identified cluster $K$
must comprise exactly the enclosed $H$ zeros throughout the slab.
On the contour, independently establish
\[
 |F|\ge m>0,\quad |H-F|\le\delta_0<m,\quad
 |(H-F)_z|\le\delta_1,\quad |F_z|\le M_1.
\]
Contour projection of the logarithmic derivative removes all enclosed
root poles and retains the exterior units. If $D_H,D_F$ are twice their
logarithmic derivatives, Note 02 proves
\begin{align}
 \sup_J|(D_H-D_F)'|&\le\kappa,\notag\\
 \kappa&=\min\left\{
 \frac{p}{\pi b^2}\frac{\delta_1+M_1\delta_0/m}{m-\delta_0},\quad
 \frac{2p}{\pi b^3}\left[-\log(1-\delta_0/m)\right]\right\}.
 \label{eq:rt-contour}
\end{align}
The scalar alternative uses the close-to-one logarithm on the contour
collar and integration by parts. Rouch\'e preserves the total zero
count with multiplicity; it supplies neither interior first-jet floors
nor individual simplicity.

In a simple-real-root regime, for the \emph{ordered} raw energy
$E_K=\sum_{i\ne j}(a_i-a_j)^{-2}$, the exterior term is
\[
 B_D=-2\sum_{i\ne j}\frac{D(a_i)-D(a_j)}{(a_i-a_j)^3},
 \qquad |B_H-B_F[a]|\le2\kappa E_K.
\]
Common motion cancels. No internal minimum gap is required for this
exterior comparison, and $B_F[a]$ is evaluated on the same true roots.
It does not assert model-root dynamics: those would add the actual
heat residual and inverse-slope costs. With fixed normalized microscopic
contour data, $p=p_0/\Ls$, $b=b_0/\Ls$ and an independent normalized
floor, Note 26 supplies $\kappa=O(\Ls^2\epsilon_X)$. A step of
length $O(\Ls^{-2})$ therefore has an $O(\epsilon_X)$ integrated
relative error coefficient. The fixed margin, geometry and their
constants remain hypotheses, not acquisitions.

This pays only the integrated error $\int2\kappa E_K\,dt$.
Cluster count/location, initial energy and the same-root model exterior
payment are still needed. For signed reference gaps $d_{ij}$, the
renormalized exterior allowance instead is
\[
 2\kappa\sum_{i\ne j}
       |(a_i-a_j)d_{ij}^{-3}-(a_i-a_j)^{-2}|,
\]
as in Note 02 (12); it is not $2\kappa\widetilde E_K$.
Raw-energy Gronwall retains an RT-sized classical reference baseline
and cannot replace the corrected additive propagation. The source's
dyadic boundary selection, reference strip, shrinking index interval
and integrated errors must all be paid together. Note 02's targeted
source-display/payment cautions retain their PDF numbering and
HTML-numbering distinction; this is not a full audit of the published
argument or its arithmetic inputs.

The exact heat pair $z^2-2s$ and $z^2-2s-\varepsilon$ has zero exterior
drift and a uniform enclosing-contour margin, while its relative ordered
energy error is $\varepsilon/(2s+\varepsilon)\to1$ as $s\downarrow0$.
Thus excellent exterior control does not establish endpoint
nonconcentration. The broader source screen has precise failed premises:
Euler-factor controls change the canonical coefficients;
weighted-square/finite-tail controls lack full reflection/completion;
reflected and positive-primitive controls lack the full theta arithmetic;
finite-prime/profile controls preserve only selected constraints;
RH-conditional endpoint controls lack ordinary Dirichlet arithmetic;
and polynomial controls lack the theta source. Generic local identities
may still apply to those controls. No control certificates were rerun.

The original negative-threshold argument also uses a real/simple-root
domain, a fixed relaxation interval and the actual zeta slice at time
zero. Translating time does not restore the last premise for a positive
threshold. No RH conclusion, new Newman upper bound, simplicity theorem
or new zero-exclusion region follows from these reviewed bounds.

\subsection{What an improved theorem should report}
A result is ready for the Book manuscript when it states: the canonical
normalization and unchanged physical data; the exact order of choices;
the full buffered domain; the positive error envelope; every retained
arithmetic and contour payment; spatial and direct time errors; the
actual model residual; and the specific quantity improved relative to
the present theorem. If it is an evaluation theorem, it must additionally
state input representation, guards, provider costs and dependence on
height, time, order and requested bits.

The current synthesis already offers a useful division: the positive
and negative recurrences explain how to improve local models, while the
source construction explains how to maintain finite global coverage.
The outstanding work is to make their interfaces equally transparent.
That is the sense in which this is a programme toward a proof from
``The Book.''

\appendix
\section{Proof sources, supporting manuscripts and status}
\label{sec:sources}

\subsection{The supporting roles}
This manuscript is the central synthesis. The existing parent manuscript
supports its physical equations, clock geometry, weak transfer,
nonholomorphic assembly and effective construction. The ansatz
manuscript supports its carriers, positive saddle, crossover and
regional bridges. The recurrence manuscript supports the negative-time
hierarchy, arithmetic completion, simultaneous order estimates and finite
providers. Their detailed proofs should remain in these technical homes;
the Book manuscript presents their common object and consequences.

As of 26 October 2026, the existing reading copies have 112, 33 and
46 pages respectively. The positive arbitrary fixed-order theorem and
compact $P_4$ bridge currently have their most complete proofs in
Ansatz Notes 19--20, not in the older 33-page ansatz manuscript.
Likewise parent Notes 58--61 extend the source/evaluation discussion
beyond the current 112-page parent manuscript. This synthesis uses those
notes directly and does not pretend the companion revisions are already
complete. The separate revision plan records their next production task.

Module numbers in the table identify source-file prefixes, not the
compiled section numbers; note numbers identify the original research notes.

\clearpage
\small
\begin{longtable}{>{\raggedright\arraybackslash}p{.24\textwidth}>{\raggedright\arraybackslash}p{.36\textwidth}>{\raggedright\arraybackslash}p{.30\textwidth}}
\toprule
Result family & Supporting proof home & Scope retained here\\
\midrule\endfirsthead
\toprule Result family & Supporting proof home & Scope retained here\\
\midrule\endhead
Canonical source, exact equations, covariance & Parent modules 01, 05; parent Notes 02, 09, 14 & Actual signed covariance; variable-regulator terms retained.\\
Negative Gaussian endpoint & Parent modules 2, 5; Notes 11, 26 & Absolute field error, adaptive physical endpoint rates.\\
Inner clock and geometry & Parent modules 3, 4, 7; Notes 12, 16, 19 & Growing inner window; exact phase and transport; unit drift retained.\\
Outer gamma and physical layer & Parent module 11; Notes 27--28; ansatz carrier sections & Named linear center and stated rounding law; prescribed switch.\\
Full saddle and leading physical transfer & Parent modules 12--13, 22; Notes 30--31, 54 & Existing physical rates, full source and buffers.\\
Compact weak and signed transfer & Parent modules 9, 15, 17; Notes 34--37 & Whole-support weak tests, actual means and source terms.\\
Finite sources and obstructions & Parent modules 8, 10, 16, 19, 24; Notes 20, 25, 37, 44, 55 & Scoped source exclusions; alternative finite constructors.\\
Spatial, moving and temporal joins & Parent modules 18, 20--23; Notes 42--50, 54 & Nonholomorphic corrections, both curvature derivatives, full tuple.\\
Source compression and evaluation & Parent Notes 58--61; Summary 62; Reviews 37--40 & Polynomial descriptor at fixed upper time; polynomial accuracy cost at fixed data/box.\\
Signed near-zero crossover & Ansatz manuscript; Notes 04, 06 & Narrow signed window and separate weaker paid join.\\
First and second positive corrections & Ansatz Notes 17--18; Reviews 12--13 & Finite coefficients, full source payment, nonzero residual.\\
Arbitrary fixed positive order & Ansatz Note 19; Review 14 & Fixed $J$ before onset; no growing-order conclusion.\\
Compact positive bridge & Ansatz Note 20; Review 15; Summary 21 & Arithmetic term retained; $P_4$ on $[63,65]$; physical transfer open.\\
Negative fixed/growing hierarchy & Recurrence modules 2--9 and underlying notes & Gamma/outer norm distinction, frozen data and exact heat evolution.\\
Negative finite providers and onset & Recurrence modules 10--12; Notes 16, 18; Reviews 12--13 & Strict guards, sharpened onset, count improvements and precision cost.\\
\bottomrule
\end{longtable}
\normalsize

\subsection{Auxiliary results and evidence}
The parent appendices contain complete-source moment bounds, derivative
allowances, clock phase formulas and finite diagnostic records. The
ansatz appendix contains exact constant-regulator cosine-cell averages,
including their potential, density and covariance. Cell averaging is
not the fixed physical convolution used in the global theorem; these
identities are diagnostic models, not a replacement closure. Earlier
small-width gradient covariance expansions carry smoothness and width
premises. The signed-source and positive Hermite studies report specific
same-mesh allowance gains and failures; sampled fits are not continuous
certificates. The recurrence component replays likewise do not replace
the full analytical bounds. This synthesis retains these results as
supporting tools and evidence rather than repeating all historical
calculations.

The proof-status distinction matters. The recent positive notes state
``proved in-note from named inputs'' with same-model internal assessment.
Other parent steps have focused cross-model LLM reviews. Neither status
equals a specialist audit of the entire inherited chain. Historical
parent tasks HDA-17 (Partial), HDA-18 (Gated) and HDA-20 (unfulfilled
full audit) are not promoted by this editorial synthesis. Bounded
completed construction/evaluation tasks retain their own narrower scope.
The new unified progress file coordinates these records without replacing
their original evidence or scientific ownership.

\subsection{Reproducibility of this edition}
The Book source is maintained in the base investigation's \texttt{book/}
folder. A hash manifest pins the current manuscript and note inputs;
the build record pins this source and its PDF. PDFs and temporary build
products live in the repository's ignored output tree. No manuscript
snapshot directory is created. The bibliography below links the principal
local proof sources, which in turn give the complete original references.
Local links assume the repository checkout path supplied at build time;
external circulation should include the companion sources or replace
those links with stable public versions before submission.

\paragraph{Acknowledgement and responsibility.}
The project and manuscript are prepared for Edward Baker. Substantial
LLM assistance contributed to drafting, derivation, checking and editorial
integration. This edition records gpt-6-astra at actual effort max;
the contributing notes retain their own model and effort records.
The author should subject the assembled work to independent mathematical
review before treating its internal assessments as publication validation.

\begin{thebibliography}{99}
\bibitem{parent} Edward Baker, with acknowledged LLM assistance,
\emph{Averaged densities for the theta heat flow: explicit asymptotics and all-time approximation}, maintained parent
manuscript, 26 October 2026 companion edition.
\src{manuscript/main.tex}{Local manuscript source}.
\bibitem{ansatz} Edward Baker, with acknowledged LLM assistance,
\emph{Evolution equations and large-height heat asymptotics: controlled ans\"atze, arithmetic modes, and physical density}, 26 October 2026 companion edition.
\src{ansatz/manuscript/main.tex}{Local manuscript source}.
\bibitem{recurrence} Edward Baker, with acknowledged LLM assistance,
\emph{Finite Fluctuation Hierarchies for the Negative-Time Riemann Heat Flow: canonical approximation, effective evaluation, and explicit bounds},
26 October 2026 working edition.
\src{recurrence/manuscript/main.tex}{Local manuscript source}.
\bibitem{positive} Edward Baker, assisted by gpt-6.1-sol, effort ultra,
\emph{Arbitrary fixed order of the unconditional full-gamma saddle},
Ansatz Note 19, 26 October 2026.
\src{ansatz/notes/19_ARBITRARY_FIXED_POSITIVE_TIME_ORDER_20261026.md}{Proof and finite constants}.
\bibitem{bridge} Edward Baker, assisted by gpt-6.1-sol, effort ultra,
\emph{First regional bridge: compact positive-time fields through the switch},
Ansatz Note 20, 26 October 2026.
\src{ansatz/notes/20_COMPACT_POSITIVE_FIELD_BRIDGE_20261026.md}{Proof and source payments}.
\bibitem{global} Edward Baker, with acknowledged LLM assistance,
\emph{Effective positive attachment and all-time existence},
parent Note 54, 25 October 2026.
\src{notes/54_EFFECTIVE_POSITIVE_ATTACHMENT_AND_ALL_TIME_EXISTENCE_20261025.md}{Complete construction}.
\bibitem{conditioning} Edward Baker, with acknowledged LLM assistance,
\emph{Fixed central conditioning and source cost}, parent Note 55,
25 October 2026.
\src{notes/55_FIXED_CENTRAL_CONDITIONING_AND_SOURCE_COST_20261025.md}{Conditioning proof}.
\bibitem{compression} Edward Baker, assisted by gpt-6.1-sol, effort ultra,
\emph{Positive analytic-source compression with a complete construction ledger},
parent Note 58, 25 October 2026.
\src{notes/58_POSITIVE_ANALYTIC_SOURCE_COMPRESSION_20261025.md}{Source construction}.
\bibitem{evaluation} Edward Baker, with acknowledged LLM assistance,
\emph{Evaluator working precision and continuous filter}, parent Note 59,
25 October 2026.
\src{notes/59_EVALUATOR_WORKING_PRECISION_AND_CONTINUOUS_FILTER_20261025.md}{Full evaluator}.
\bibitem{phase} Edward Baker, with acknowledged LLM assistance,
\emph{Exact phase index by dyadic onset selection}, parent Note 60,
25 October 2026.
\src{notes/60_EXACT_PHASE_INDEX_BY_DYADIC_ONSET_SELECTION_20261025.md}{Phase selection}.
\bibitem{cost} Edward Baker, assisted by gpt-6.1-sol, effort ultra,
\emph{Polynomial accuracy-bit cost for the complete evaluator and filter},
parent Note 61, 25 October 2026.
\src{notes/61_POLYNOMIAL_ACCURACY_EVALUATOR_AND_FILTER_COST_20261025.md}{Cost theorem and weak-primitive frontier}.
\bibitem{rec-onset} Edward Baker, with acknowledged LLM assistance,
\emph{Sharper all-height onset}, recurrence Note 16, 26 October 2026.
\src{recurrence/notes/16_SHARPER_ALL_HEIGHT_ONSET_20261026.md}{Onset theorem}.
\bibitem{rec-bounds} Edward Baker, with acknowledged LLM assistance,
\emph{Uniform Riemann--Siegel bound strengthening}, recurrence Note 18,
26 October 2026.
\src{recurrence/notes/18_UNIFORM_RS_BOUND_STRENGTHENING_20261026.md}{Uniform bounds and count savings}.
\bibitem{centers} Edward Baker, assisted by gpt-6.1-sol, effort ultra,
\emph{Named centers at the physical regulator switch}, ansatz Note 22,
26 October 2026, closed A03-AN01 return, same-model internal assessment.
\src{ansatz/notes/22_NAMED_SWITCH_CENTERS_20261026.md}{Original proof note}; \src{ansatz/reviews/16_NAMED_SWITCH_CENTERS_ASSESSMENT_20261026.md}{same-model assessment}.
\bibitem{rec-moving} Edward Baker, assisted by gpt-6.1-sol, inherited effort ultra,
\emph{The specified moving lower endpoint $\tau(L)=1/L$}, recurrence Note 20,
26 October 2026, closed R05-AN08 return, same-model internal assessment.
\src{recurrence/notes/20_MOVING_LOWER_ENDPOINT_20261026.md}{Original proof note and endpoint interface}; \src{recurrence/reviews/14_MOVING_ENDPOINT_ASSESSMENT_20261026.md}{same-model assessment}.
\bibitem{p4-physical} Edward Baker, assisted by gpt-6.1-sol, effort ultra,
\emph{Fixed-$P_4$ physical tuple on the compact regulator switch},
Ansatz Note 23 and Assessment 17, 26 October 2026.
\src{ansatz/notes/23_FIXED_P4_PHYSICAL_TUPLE_20261026.md}{Proof provenance};
\src{ansatz/reviews/17_FIXED_P4_PHYSICAL_TUPLE_ASSESSMENT_20261026.md}{same-model assessment}.
\bibitem{small-positive} Edward Baker, assisted by gpt-6.1-sol, effort ultra,
\emph{First small-positive slab: a complete buffered canonical comparison},
Ansatz Note 24 and Assessment 18, 26 October 2026.
\src{ansatz/notes/24_SMALL_POSITIVE_SLAB_20261026.md}{Proof provenance};
\src{ansatz/reviews/18_SMALL_POSITIVE_SLAB_ASSESSMENT_20261026.md}{same-model assessment}.
\bibitem{intermediate} Edward Baker, assisted by gpt-6.1-sol, effort ultra,
\emph{Corrected buffered comparison on $[3/4,33]$ and its two overlaps},
Ansatz Note 25 and Assessment 19, 26 October 2026.
\src{ansatz/notes/25_INTERMEDIATE_POSITIVE_COMPARISON_20261026.md}{Proof provenance};
\src{ansatz/reviews/19_INTERMEDIATE_POSITIVE_COMPARISON_ASSESSMENT_20261026.md}{same-model assessment}.
\bibitem{retained} Edward Baker, assisted by gpt-6.1-sol, effort ultra,
\emph{Retained arithmetic on the moving positive band and both overlaps},
Ansatz Note 26 and Assessment 20, 26 October 2026.
\src{ansatz/notes/26_RETAINED_ARITHMETIC_MOVING_POSITIVE_BAND_20261026.md}{Proof provenance};
\src{ansatz/reviews/20_RETAINED_ARITHMETIC_MOVING_BAND_ASSESSMENT_20261026.md}{same-model assessment}.
\bibitem{rt-interface} Edward Baker, assisted by gpt-6.1-sol, effort ultra,
\emph{Corrected propagation and a contour requirement for exterior force},
Rodgers--Tao branch Note 02 and Assessment 01, 26 October 2026.
\src{../02_heat_flow_rh_applications/rodgers_tao/notes/02_CORRECTED_PROPAGATION_AND_CONTOUR_FORCE_TRANSFER_20261026.md}{Conditional interface and bounded source audit};
\src{../02_heat_flow_rh_applications/rodgers_tao/reviews/01_PROPAGATION_AND_CONTOUR_TRANSFER_ASSESSMENT_20261026.md}{same-model assessment}.
The note retains PDF numbering and identifies the different HTML numbering.
Its sole \src{../02_heat_flow_rh_applications/rodgers_tao/research-progress/01_RODGERS_TAO_RESEARCH_PROGRESS_20261026.md}{application register}
owns RT-03/RT-04 r2 and later work.
\bibitem{poly} D. H. J. Polymath,
\emph{Effective approximation of heat flow evolution of the Riemann
$\xi$ function, and a new upper bound for the de Bruijn--Newman constant},
arXiv:1904.12438v2 (2019), especially equation 35 and Proposition 6.2.
\href{https://arxiv.org/abs/1904.12438v2}{Primary reference}.
The precise uses and normalizations are recorded in the supporting proofs.
\bibitem{dlmf} NIST Digital Library of Mathematical Functions,
\href{https://dlmf.nist.gov/5.11}{Section 5.11, gamma asymptotics and error bounds},
and \href{https://dlmf.nist.gov/24.8}{Section 24.8, Bernoulli Fourier series}.
Consulted 26 October 2026; the supporting notes specify their sector,
branch and derivative uses.
\bibitem{weak-consumer} Edward Baker, assisted by gpt-6.1-sol, effort ultra,
\emph{Complete weak-filter puncture consumer and boundary moduli},
Note 71 and Review 46, 26 October 2026.
\src{notes/71_WEAK_FILTER_PUNCTURE_CONSUMER_20261026.md}{Proof and finite-data contract}.
\bibitem{graded-cost} Edward Baker, assisted by gpt-6.1-sol, effort ultra,
\emph{Graded weak filtering and complete acquisition cost},
Note 72 and Review 47, 26 October 2026.
\src{notes/72_GRADED_WEAK_FILTER_AND_ACQUISITION_COST_20261026.md}{Fixed-family proof}.
\bibitem{binary-p} Edward Baker, assisted by gpt-6.1-sol, effort ultra,
\emph{Binary-height cost for the complete pure-P exterior query},
Note 73 and Review 48, 26 October 2026.
\src{notes/73_BINARY_HEIGHT_PURE_P_EXTERIOR_COST_20261026.md}{Exterior query proof}.
\bibitem{all-height-cost} Edward Baker, assisted by gpt-6.1-sol, effort ultra,
\emph{Complete all-height query-cost ledger on $[95,97]$},
Note 74 and Review 49, 26 October 2026.
\src{notes/74_ALL_HEIGHT_QUERY_COST_LEDGER_20261026.md}{Complete cover and constructor ledger}.
\bibitem{binary-g} Edward Baker, assisted by gpt-6.1-sol, effort ultra,
\emph{Complete binary-height pure-G query cost on $[62,63]$},
Note 75 and Review 50, 26 October 2026.
\src{notes/75_BINARY_HEIGHT_PURE_G_QUERY_COST_20261026.md}{Endpoint provider and fixed-floor proof}.
\bibitem{summary76} Edward Baker, assisted by gpt-6.1-sol, effort ultra,
\emph{Chat0028 research summary}, Note 76, 26 October 2026.
\src{notes/76_CHAT0028_RESEARCH_SUMMARY_20261026.md}{Five-stage result and scope map}.
\bibitem{rt-pilot} Edward Baker, assisted by gpt-6.1-sol, effort ultra,
\emph{Canonical two-root margins and complete force budget},
Rodgers--Tao Note 03 and Assessment 02, 26 October 2026.
\src{../02_heat_flow_rh_applications/rodgers_tao/notes/03_CANONICAL_TWO_ROOT_MARGINS_AND_FORCE_BUDGET_20261026.md}{Finite certificate and comparison}.
\bibitem{simplicity-targets} Edward Baker, assisted by gpt-6-astra, effort max,
\emph{Target bounds for RH-conditional simplicity}, Note 02 and
Assessment 02, 26 October 2026.
\src{../02_heat_flow_rh_applications/rh_simple_zeros/notes/02_TARGET_BOUNDS_FOR_RH_TO_SIMPLICITY_20261026.md}{J/P/V/E definitions and derived implications}.
\bibitem{target-review} Edward Baker, assisted by gpt-6-astra, effort max,
\emph{Heat-density research against the RH and simplicity targets},
Review 51, 26 October 2026.
\src{reviews/51_RH_TARGET_GAP_ANALYSIS_AND_RESEARCH_PRIORITIES_20261026.md}{Source-grounded assessment and research priorities}.
\end{thebibliography}

\end{document}
